\documentclass[a4paper,11pt]{article}

\usepackage{amsmath}
\usepackage{array}
\usepackage{amsthm}
\usepackage{ifthen}
\usepackage[utf8]{inputenc}
\usepackage[adobe-utopia]{mathdesign}
\usepackage[T1]{fontenc}
\usepackage{graphicx}
\usepackage[top=2.5cm,bottom=2.5cm,left=2cm,right=2cm]{geometry}
%\setlength{\parskip}{\baselineskip}
\usepackage[spanish,es-noquoting]{babel}
\usepackage{float}
\usepackage[    pdftex,unicode=true,
        pdfdisplaydoctitle=true,
        colorlinks,
        linkcolor=black,]{hyperref}

\usepackage[usenames,dvipsnames]{xcolor}

\usepackage[font={scriptsize,sf},format=plain,labelfont={bf},labelsep=period,skip=-0.5pt,justification=centering]{caption}
%\usepackage[font+=scriptsize,skip=-0.5pt]{subcaption}

%\usepackage{indentfirst}
%\usepackage{placeins}
\usepackage{tikz}
\usetikzlibrary{decorations.pathreplacing}
\usetikzlibrary{positioning}
\usetikzlibrary{calc}
\usetikzlibrary{shapes.geometric}
\usetikzlibrary{babel}
\usepackage{pgfplots}
\pgfplotsset{compat=1.18}
\usepackage{circuitikz}
%\usepackage{rotating}
%\usepackage{lscape}
%\usepackage{pifont}
\usepackage{siunitx}
\DeclareSIUnit{\baud}{baudios}
\DeclareMathOperator{\sinc}{sinc}
%environment problema
\newcounter{numejer}[section]
\newenvironment{ejercicio}[1][]{
\addtocounter{numejer}{1} 
\noindent \large \textbf{Ejercicio \arabic{section}.\arabic{numejer}} \ifthenelse{\equal{#1}{}}{}{\emph{(#1)}}

\smallskip
}{
\bigskip
}

\newenvironment{solucion}[1][]{
\noindent \large \textsl{Solución:} \ifthenelse{\equal{#1}{}}{}{\emph{(#1)}}

\smallskip
}{
\bigskip
}

\newenvironment{blockdiagram}[1][]{%
\begin{tikzpicture}[
     block/.style={
    % The shape:
    rectangle,
    % The size:
    minimum size=6mm,draw,
    % The border:
    thick,
    },
    operator/.style={
    circle,draw,
    minimum size=7mm,
    inner sep=1mm,
    thick},
    punto/.style={coordinate},auto,
    etiqueta/.style={rectangle,inner sep=0mm},
    multD/.style={isosceles triangle,draw,isosceles triangle apex angle=50,shape border rotate =-90,thick},#1]
}{\end{tikzpicture}}


\title{Prácticos de Modulación Digital -- Soluciones}
\author{Universidad ORT Uruguay}
\date{Curso 2024}

\begin{document}
\maketitle

\section{Modulación PAM en banda base}

\begin{ejercicio}
Un sistema de comunicación digital envía una señal PAM digital NRZ cuaternaria, con símbolos independientes $a_k=\{0,A,2A,3A\}$. Las probabilidades de ocurrencia de cada símbolo son $2/5$, $2/5$, $1/10$ y $1/10$, respectivamente. Calcule y bosqueje la densidad espectral de potencia de la señal trasmitida.
\end{ejercicio}

\begin{solucion}

Calculamos primero $\mu$ y $\sigma^2$ de los símbolos:
\begin{align*}
        \mu &= 0 (2/5) + A (2/5) + 2A (1/10) + 3A (1/10) = 0.9 A \Rightarrow \mu^2 = 0.81 A^2\\
        E[a_k^2] &= 0 (2/5) + A^2 (2/5) + 4A^2 (1/10) + 9A^2 (1/10) = 1.7 A^2\\
        \sigma^2 &= E[a_k^2] - \mu^2 = 0.89 A^2
\end{align*}

Recordamos ahora la fórmula de la densidad espectral de potencia:
\begin{equation*}
        S_x(f) = \sigma^2 f_s |\hat{P}(f)|^2 + \mu^2 f_s^2 \sum_k |P(kf_s)|^2 \delta(f-k f_s)
\end{equation*}

Como $p(t) = \mathbf{1}_{[0,T_s)}(t)$ se tiene que:
\begin{equation*}
        |\hat{P}(f)| = T_s \sinc (\pi T_s f) \; \Rightarrow \; |\hat{P}(f)|^2 = T_s^2 \sinc^2 (\pi T_s f)
\end{equation*}
Observemos que el $\sinc$ se anula en $k f_s$ por lo que solo hay término de continua en la parte discreta. El espectro final es:
\begin{equation*}
        S_x(f) = 0.89 A^2 T_s \sinc^2(\pi T_s f) + 0.81 A^2 \delta(0)
\end{equation*}
\begin{center}
\begin{tikzpicture}
        \pgfmathdeclarefunction{sinc}{1}{%
        \pgfmathparse{(#1==0 ? 1: sin(#1 r))/(#1==0 ? 1: #1)}%
        }
        \begin{axis}[x=2cm, y=2cm, axis lines=middle,  
                xlabel=$f$,
                x label style={anchor=north},
                y label style={anchor=east},
                xmin=-3,xmax=3,ymin=0,ymax=1.5,
                xtick={-2,-1,0,1,2}, ytick=\empty,
                xticklabels={$-2f_s$,$-f_s$, $0$, $f_s$, $2f_s$},
                clip=false,
                smooth
            ]
        \addplot+[thick,mark=none,teal, domain = (-3:3), samples=100] {.89*sinc(pi*x)^2} node[midway,above right] {$0.89 A^2 T_s$};

        \draw[-stealth,very thick, teal] (axis cs:0,0) -- (axis cs:0,.81) node[below right] {$0.81 A^2$};
        \end{axis}
\end{tikzpicture}
\end{center}

\end{solucion}

\begin{ejercicio}
        Sea una señal PAM digital polar binaria $x(t)$ que utiliza pulsos rectangulares $p(t)$ de ancho $T_s$, con formato RZ (retorno a cero en $T_z<T_s$).
        \begin{enumerate}
        \item Asumiendo que los símbolos $a_k=\{A,-A\}$ son independientes y equiprobables, bosqueje la densidad espectral de potencia y calcula la potencia de la señal $\langle x^2(t)\rangle$.
        \item Supongamos ahora que la trasmisión es en formato NRZ y $P(A)=\alpha$ y $P(-A)=1-\alpha$, es decir, los símbolos no son equiprobables. Demuestre que $\langle x^2(t)\rangle$ no depende de $\alpha$.
        \end{enumerate}                
\end{ejercicio}
        
\begin{solucion}

        \begin{enumerate}
        \item Calculamos primero $\mu$ y $\sigma^2$ de los símbolos:
        \begin{align*}
                \mu &= -A (1/2) + A (1/2) = 0 \\
                \sigma^2 &= E[a_k^2] = A^2 (1/2) + A^2(1/2) = A^2
        \end{align*}
        
        Calculamos ahora la transformada del pulso de tiempo $T_z<T_s$:
        \begin{equation*}
                |\hat{P}(f)| = T_z \sinc (\pi T_z f) \; \Rightarrow \; |\hat{P}(f)|^2 = T_z^2 \sinc^2 (\pi T_z f)
        \end{equation*}
        Donde se usó que el pulso es de ancho $T_z$.

        Como $\mu=0$ la fórmula de densidad espectral de potencia queda:
        \begin{equation*}
                S_x(f) = \sigma^2 f_s |\hat{P}(f)|^2
        \end{equation*}
        
        El espectro final es:
        \begin{equation*}
                S_x(f) = A^2 f_s T_z^2 \sinc^2(\pi T_z f)
        \end{equation*}
        Notar que $f_s$ y $f_z = 1/T_z$ ahora son distintos!! Los pulsos son más rápidos que el tiempo de símbolo ($T_z<T_s \Rightarrow f_z > f_s$) por lo que ocupa más ancho de banda.

        Ejemplo: $T_z = T_s/2$, medio símbolo, entonces $f_z = 2f_s$.
        \begin{center}
        \begin{tikzpicture}
                \pgfmathdeclarefunction{sinc}{1}{%
                \pgfmathparse{(#1==0 ? 1: sin(#1 r))/(#1==0 ? 1: #1)}%
                }
                \begin{axis}[x=2cm, y=2cm, axis lines=middle,  
                        xlabel=$f$,
                        x label style={anchor=north},
                        y label style={anchor=east},
                        xmin=-3,xmax=3,ymin=0,ymax=1.5,
                        xtick={-2,-1,0,1,2}, ytick=\empty,
                        xticklabels={$-2f_s$,$-f_s$, $0$, $f_s$, $2f_s$},
                        clip=false,
                        smooth
                    ]
                \addplot+[thick,mark=none,teal, domain = (-3:3), samples=100] {sinc(pi*.5*x)^2} node[midway,above right] {$\frac{A^2 T_s}{4}$};
        
                \end{axis}
        \end{tikzpicture}
        \end{center}
        
        Para hallar la potencia total hay que integrar $S_x(f)$ de la siguiente forma (usando Parseval):
        \begin{equation*}
                S = \int_{-\infty}^\infty S_x(f)df = \sigma^2 f_s \int_{-\infty}^\infty |\hat{P}(f)|^2 df = \sigma^2 f_s \int_{-\infty}^\infty p^2(t)dt,
        \end{equation*}
        De donde:
        \begin{equation*}
                S = \sigma^2 f_s T_z = A^2 f_s T_z.
        \end{equation*}

        
        \item Si ahora $P(A) = \alpha$, los símbolos no son equiprobables, recalculamos $\mu$ y $\sigma^2$
        
        \begin{align*}
                \mu &= -A (1-\alpha) + A (\alpha) = A(2\alpha-1) \\
                E[a_k^2] &= A^2 (1-\alpha) + A^2(\alpha) = A^2 \\
                \sigma^2 &= A^2 - A^2(2\alpha-1)^2 = 4\alpha(1-\alpha)A^2.
        \end{align*}

        El espectro ahora presenta parte discreta, ya que $\mu\neq 0$ y $\hat{P}(f)$ no necesariamente se anula en los múltiplos de $f_s$ (se anula en múltiplos de $f_z$).

        El espectro completo es:
        \begin{equation*}
                S_x(f) = 4\alpha(1-\alpha)A^2 f_s T_z^2 \sinc^2(\pi T_z f) + A^2(2\alpha-1)^2 \sum_{k} T_z^2 \sinc^2(\pi T_z k f_s) \delta(f-kf_s)
        \end{equation*}
        Integrando en todas las frecuencias y usando que la parte continua no cambia su integral:
        \begin{equation*}
                S = \int_{-\infty}^\infty S_x(f)df = 4\alpha(1-\alpha)A^2 f_s T_z + A^2(2\alpha-1)^2 f_s  \sum_k \left[\frac{T_z}{T_s} \sinc(k\pi T_z f_s)\right]^2.
        \end{equation*}

        Notar que el último término tiene forma de Parseval de una serie de Fourier. De hecho, los coeficientes $c_k = \frac{T_z}{T_s} \sinc(k\pi T_z f_s)$ corresponden a una onda cuadrada de ancho $T_z$, período $T_s$ y altura $1$. Usando Parseval, la serie de coeficientes es entonces igual a la integral en un período de dicho pulso al cuadrado, es decir $T_z$.

        Sustituyendo queda:
        \begin{equation*}
                S = 4\alpha(1-\alpha)A^2 f_s T_z + A^2(2\alpha-1)^2 f_s T_z = A^2 f_s T_z
        \end{equation*}
        que no depende de $\alpha$ (en particular es el mismo valor que antes, $\alpha=1/2$).

        La interpretación obvia del resultado es que independientemente de $\alpha$, $<x^2>$ es una onda cuadrada de alto $A^2$ y ancho $T_z$ y período $f_s$, por lo que la potencia media transmitida es $A^2T_z f_s$, el valor medio de dicha onda.
        \end{enumerate}
\end{solucion}
        


\begin{ejercicio}
Una computadora genera palabras de 16 bits a una velocidad de 20.000 palabras por segundo. Determinar el ancho de banda mínimo necesario para trasmitirlas usando un sistema de señalización binario. Calcule el número de símbolos mínimo de un sistema M-ario si el ancho de banda de trasmisión disponible es de 60\unit{\kHz}.
\end{ejercicio}

\begin{solucion}
        Velocidad de transmisión de información requerida:
        \begin{equation*}
                \mathrm{vti} = 16 \times 20000 = 320 \, \textrm{kbps}
        \end{equation*}
        El sistema es binario, por lo que el mínimo ancho de banda corresponde a usar pulsos de Nyquist con $\mathrm{vti} = f_s = 2B$ de donde el mínimo es $160 \unit{\kilo\hertz}$.

        Si en cambio, el ancho de banda disponible es $60 \unit{\kilo\hertz}$, la frecuencia de símbolos máxima es $f_s=2B=120 \textrm{kbaud/s}$. Necesito hallar $M$ tal que:
        \begin{equation*}
                 \mathrm{vti} = \log_2(M) f_s \Rightarrow M = 2^{\frac{320}{120}} \approx 6.35 \mapsto 7 \textrm{ bits}
        \end{equation*}
        En la práctica usaríamos $8=2^3$ y un ancho de banda algo menor ($\approx 53 \unit{\kilo\hertz}$).
\end{solucion}

\begin{ejercicio}
        Un dispositivo debe transmitir información a una tasa de 56kbps utilizando señalización binaria y pulsos de Nyquist con exceso de ancho de banda $\alpha$.
        \begin{enumerate}
        \item Calcule el ancho de banda necesario para la transmisión para $\alpha\in\left\{0,\frac{1}{4},\frac{1}{2},1\right\}$.
        \item Suponga ahora que el canal tiene un ancho de banda de 30kHz y se utiliza señalización $M$-aria. Calcule el mínimo valor de $M$ necesario en cada caso para los valores de $\alpha$ utilizados en la parte anterior.
        \end{enumerate}
\end{ejercicio}

\begin{solucion}

        \begin{enumerate}
        \item $\textrm{vti} = 56 \textrm{ kbps}$, señalización binaria por lo que $f_s=\textrm{vti} = 56 \unit{\kilo\hertz}$. Para pulsos de Nyquist con exceso $\alpha$ vale:
        \begin{equation*}
                \frac{f_s}{2}(1+\alpha) = B
        \end{equation*}
        De donde los anchos de banda son $28 \unit{\kilo\hertz}$, $35 \unit{\kilo\hertz}$, $42 \unit{\kilo\hertz}$ y $56 \unit{\kilo\hertz}$ respectivamente.

        \item Si el canal es de $30 \unit{\kilo\hertz}$, las posibles frecuencias de símbolo son $f_s = \frac{2B}{1+\alpha}$, es decir $f_s = 60 \textrm{ kbaud}$, $48 \textrm{ kbaud}$, $40 \textrm{ kbaud}$ y $30 \textrm{ kbaud}$ respectivamente. Los $M$ necesarios son $2^{\frac{\textrm{vti}}{f_s}}$ lo que da $M=2$ para $\alpha=0$, $M=3$ para los intermedios y $M=4$ para $\alpha=1$ (redondeando al entero más cercano).
        \end{enumerate}
\end{solucion}

\begin{ejercicio}
En la figura se ilustran los espectros de dos pulsos, $\hat{g}(f)$ y $\hat{h}(f)$, con $0<\alpha<1$.
\begin{figure}[ht]
\centering
\begin{tikzpicture}
\begin{axis}[   x=3cm, y=1.5cm, axis lines=middle,  
                xlabel=$f$,
                x label style={anchor=north},
                y label style={anchor=east},
                xmin=-1,xmax=1,ymin=-.2,ymax=2,
                xtick={-.1,.1}, ytick=\empty,
                xticklabels={$\scriptstyle -\frac{\alpha}{2T_s}$,$\scriptstyle \frac{\alpha}{2T_s}$},
                clip=false,
                smooth
            ]
\addplot+[thick,mark=none,const plot, teal] coordinates
{(-.9,0) (-.1,0) (-.1,1.4) (.1,1.4) (.1,0) (.9,0)} node[above left,midway]{\small $T_s/\alpha$};
\node at (axis cs:-.8,1.5) {$\hat{g}(f)$};
\end{axis}
\end{tikzpicture}\hspace{1cm}
\begin{tikzpicture}
\begin{axis}[   x=3cm, y=1.5cm, axis lines=middle,  
        xlabel=$f$,
        x label style={anchor=north},
        y label style={anchor=east},
        xmin=-1,xmax=1,ymin=-.2,ymax=2,
        xtick={-.8,.8}, ytick=\empty,
        xticklabels={$\scriptstyle -\frac{1}{2T_s}$,$\scriptstyle \frac{1}{2T_s}$},
        clip=false,
        smooth
    ]
\addplot+[thick,mark=none,const plot, teal] coordinates
{(-.9,0) (-.8,0) (-.8,1) (.8,1) (.8,0) (.9,0)} node[above left,midway]{\small $T_s$};
\node at (axis cs:-.8,1.5) {$\hat{h}(f)$};
\end{axis}
\end{tikzpicture}
\end{figure}

\begin{enumerate}
\item Bosquejar el espectro $\hat{p}(f)=(\hat{g}*\hat{h})(f)$.
\item Pruebe que $p(t)=\mathcal{F}^{-1}\{\hat{p}(f)\}$ cumple con la condición de ISI nula y calcule $p(0)$.
\item El pulso $p(t)$ se utiliza para generar una señal PAM digital $x(t) = \sum_k a_kp(t-kT_s),$
donde $a_k=\{-3A,-A,A,3A\}$ es una secuencia de símbolos
independientes y equiprobables. Exprese analíticamente la densidad
espectral de potencia de la señal $x(t)$.
\item En las condiciones de la parte anterior y asumiendo $\alpha=\frac{2}{3}$,
calcule la relación entre la velocidad de trasmisión de información y el ancho de banda del canal de trasmisión.
\end{enumerate}
\end{ejercicio}

\begin{solucion}
\begin{enumerate}

        \item La convolución de estos pulsos, al ser $0<\alpha<1$ produce el siguiente trapecio:

        \begin{center}
                \begin{tikzpicture}
                        \begin{axis}[   x=5cm, y=1.5cm, axis lines=middle,  
                                xlabel=$f$,
                                x label style={anchor=north},
                                y label style={anchor=east},
                                xmin=-1.2,xmax=1.2,ymin=-.2,ymax=2,
                                xtick={-1,-.85,-.7,0,.7,.85,1}, ytick=\empty,
                                xticklabels={$\scriptstyle -\frac{1+\alpha}{2T_s}$, $\scriptstyle -\frac{1}{2T_s}$, $\scriptstyle -\frac{1-\alpha}{2T_s}$, 0 , $\scriptstyle \frac{1-\alpha}{2T_s}$, $\scriptstyle \frac{1}{2T_s}$, $\scriptstyle \frac{1+\alpha}{2T_s}$},
                                clip=false,
                            ]
                        \addplot[very thick,mark=none, teal] coordinates
                        {(-1.2,0) (-1,0) (-.7,1) (.7,1) (1,0) (1.2,0)} node[above left,midway]{$T_s$};
                        \node at (axis cs:-.8,1.5) {$\hat{p}(f)$};
                        \end{axis}
                \end{tikzpicture}                        
        \end{center}

        \item Podemos hacerlo de dos maneras:
         \begin{itemize}
          \item En frecuencia, ver que tiene simetría vestigial alrededor de $1/(2T_s)$ por lo que usando $f_s= 1/T_s$ no debería introducir ISI. $p(0)$ es el área bajo la gráfica de $\hat{p}(f)$, $p(0)=1$.
          \item En el tiempo, observando que la antitransformada del pulso es el producto de dos $\sinc$:
           \begin{equation*}
                p(t) = \sinc(\pi f_s t)\sinc(\pi \alpha f_s t).
           \end{equation*}
           De aquí se deduce que $p(0)=1$ y que, por el primer término, no hay ISI a frecuencia $f_s=1/T_s$.
         \end{itemize}

        \item Observemos que $\mu=0$ y $\sigma^2 = E[a_k^2] = (-3A)^2(1/4) + (-A)^2(1/4) + A^2 (1/4) + (3A)^2 (1/4) = 5A^2$. Por lo tanto el espectro es solo continuo según la fórmula:
         \begin{equation*}
                S_x(f) = \sigma^2 f_s |\hat{p}(f)|^2.
         \end{equation*}
         Se obtiene por lo tanto:
         \begin{equation*}
                S_x(f) = \begin{cases}
                5A^2 T_s^3 \left(f+\frac{1+\alpha}{2T_s}\right)^2 \frac{1}{\alpha^2} & -\frac{1+\alpha}{2T_s} \leqslant f \leqslant -\frac{1-\alpha}{2T_s} \\
                5A^2 T_s & -\frac{1-\alpha}{2T_s} \leqslant f \leqslant \frac{1-\alpha}{2T_s} \\
                5A^2 T_s^3 \left(f-\frac{1+\alpha}{2T_s}\right)^2 \frac{1}{\alpha^2}& \frac{1-\alpha}{2T_s} \leqslant f \leqslant \frac{1+\alpha}{2T_s} \\
                0 & \textrm{fuera de la banda}.
                \end{cases}
         \end{equation*}
         Para hallar las regiones laterales se elevó al cuadrado la ecuación de la recta del trapecio correspondiente.

         Bosquejo:
         \begin{center}
                \begin{tikzpicture}
                        \begin{axis}[   x=5cm, y=1.5cm, axis lines=middle,  
                                xlabel=$f$,
                                ylabel=$S_x(f)$,
                                x label style={anchor=north},
                                y label style={anchor=east},
                                xmin=-1.2,xmax=1.2,ymin=-.2,ymax=2,
                                xtick={-1,-.85,-.7,0,.7,.85,1}, ytick=\empty,
                                xticklabels={$\scriptstyle -\frac{1+\alpha}{2T_s}$, $\scriptstyle -\frac{1}{2T_s}$, $\scriptstyle -\frac{1-\alpha}{2T_s}$, 0 , $\scriptstyle \frac{1-\alpha}{2T_s}$, $\scriptstyle \frac{1}{2T_s}$, $\scriptstyle \frac{1+\alpha}{2T_s}$},
                                clip=false,
                            ]
                        \addplot+[very thick,mark=none, const plot, teal, domain=(-.7:.7)] {1} node[above left,midway]{$5A^2T_s$};
                        \addplot+[thick,mark=none, teal, domain=(-1:-.7)] {(x+1)^2/.09};
                        \addplot+[thick,mark=none, teal, domain=(.7:1)] {(1-x)^2/.09};
                        \end{axis}
                \end{tikzpicture}                        
        \end{center}

        \item Si $\alpha = 2/3$, el ancho de banda para que el pulso pase entero sin distorsión (e introduzca ISI a la salida) es $B = \frac{(1+2/3)}{2T_s} = \frac{5}{6T_s} = \frac{5}{6}f_s$. La velocidad de transmisión de información es $\textrm{vti} = \log_2{4} f_s = 2 f_s$, de donde la eficiencia espectral es:
        \begin{equation*}
                \textrm{Eff} = \frac{\textrm{vti}}{B} = \frac{2 f_s}{(5/6) f_s} = 2.4 \unit{\bit/\hertz}
        \end{equation*}
\end{enumerate}
\end{solucion}

\begin{ejercicio}
Se considera el pulso de duración $T$, dado en $|t|\leqslant \frac{T}{2}$ por $ p(t)=1 - \frac{2 |t|}{T}$.  Se utiliza este pulso para una transmisión digital $\sum_k a_k p(t-kT_s)$, con símbolos unipolares binarios $a_k \in \{0,A\}$.

\begin{enumerate}
\item Supongamos primero que $T_s=T$.  Bosquejar la forma de onda para el mensaje binario  1-0-1-1-1-1-0. Calcular  y bosquejar la densidad espectral de potencia suponiendo símbolos independientes y equiprobables.
\item Repetir las preguntas de la parte anterior siendo ahora $T_s=\frac{T}{2}$.
\end{enumerate}
\end{ejercicio}

\begin{solucion}
Antes de comenzar el ejercicio, analicemos el pulso en cuestión con $T$ genérico. Bosquejamos:
\begin{center}
        \begin{tikzpicture}
                \begin{axis}[   x=3cm, y=1.5cm, axis lines=middle,  
                        xlabel=$t$,
                        ylabel=$p(t)$,
                        x label style={anchor=north},
                        y label style={anchor=east},
                        xmin=-1.2,xmax=1.2,ymin=-.2,ymax=1.4,
                        xtick={-.3,0,.3}, ytick={1},
                        xticklabels={$-\frac{T}{2}$, 0 , $\frac{T}{2}$},
                        clip=false,
                    ]
                \addplot+[very thick,mark=none, teal] coordinates {(-1,0) (-.3,0) (0,1) (.3,0) (1,0)};
                \end{axis}
        \end{tikzpicture}                        
\end{center}

Observación: $p(t)$ puede escribirse como la convolución de dos pulsos de ancho $T/2$ y alto $\sqrt{2/T}$. De aquí se deduce que la transformada de Fourier es
\begin{equation*}
        \hat{P}(f) = \frac{T}{2}\sinc^2\left(\pi \frac{T}{2} f\right)
\end{equation*}

\begin{enumerate}

        \item En primer lugar, el ancho de pulso corresponde al tiempo de símbolo, por lo que la onda PAM requerida queda:
        \begin{center}
                \begin{tikzpicture}
                        \begin{axis}[   x=1.8cm, y=1.5cm, axis lines=middle,  
                                xlabel=$t$,
                                ylabel=$x(t)$,
                                x label style={anchor=north},
                                y label style={anchor=east},
                                xmin=-1,xmax=7,ymin=-.2,ymax=1.4,
                                xtick={0,1,2,3,4,5,6}, ytick={1},
                                xticklabels={$0$,$T_s$,$2T_s$,$3T_s$,$4T_s$,$5T_s$,$6T_s$},
                                yticklabels={$A$},
                                clip=false,
                            ]
                        \addplot+[mark=none, dashed, red!50!white] coordinates {(-.5,0) (0,1) (.5,0)};
                        \addplot+[mark=none, dashed, red!50!white] coordinates {(.5,0) (1,1) (1.5,0)};
                        \addplot+[mark=none, dashed, red!50!white] coordinates {(1.5,0) (2,1) (2.5,0)};
                        \addplot+[mark=none, dashed, red!50!white] coordinates {(2.5,0) (3,1) (3.5,0)};
                        \addplot+[mark=none, dashed, red!50!white] coordinates {(3.5,0) (4,1) (4.5,0)};
                        \addplot+[mark=none, dashed, red!50!white] coordinates {(4.5,0) (5,1) (5.5,0)};
                        \addplot+[mark=none, dashed, red!50!white] coordinates {(5.5,0) (6,1) (6.5,0)};

                        \addplot+[mark=none, solid, very thick, teal] coordinates {(-.5,0) (0,1) (0.5,0) (1.5,0) (2,1) (2.5,0) (3,1) (3.5,0) (4,1) (4.5,0) (5,1) (5.5,0) (6.5,0)};
                        \end{axis}
                \end{tikzpicture}                        
        \end{center}
        Notar que el pulso dura menos de un tiempo de símbolo por lo cual no habrá ISI.

        Para calcular la densidad espectral de potencia, observamos que $\mu = A/2$, $E[a_k^2] = A^2/2$ y por lo tanto $\mu^2=A^2/4$ y $\sigma^2=A^2/4$. Además, $|\hat{P}(f)|^2 = \frac{T_s^2}{4}\sinc^4\left(\pi\frac{T_s}{2}f\right)$. De esto se deduce que:
        \begin{itemize}
                \item El espectro tiene una parte continua, que se anula en los múltiplos de $2f_s$.
                \item La parte discreta incluye un valor de continua, y deltas en todos los múltiplos impares de $f_s$ (similar al pulso cuadrado RZ con ancho medio tiempo de símbolo de las notas).
        \end{itemize}
        La expresión completa es:
        \begin{equation*}
                S_x(f) = \frac{A^2}{16}T_s \sinc^4\left(\pi\frac{f}{2f_s}\right) + \frac{A^2}{16} \sum_k \sinc^4\left(\pi \frac{k}{2}\right) \delta(f-kf_s)
        \end{equation*}

        \begin{center}
                \begin{tikzpicture}
                        \pgfmathdeclarefunction{sinc}{1}{%
                        \pgfmathparse{(#1==0 ? 1: sin(#1 r))/(#1==0 ? 1: #1)}%
                        }
                        \begin{axis}[   x=1.8cm, y=1.5cm, axis lines=middle,  
                                xlabel=$f$,
                                ylabel=$S_x(f)$,
                                x label style={anchor=north},
                                y label style={anchor=east},
                                xmin=-3.5,xmax=3.5,ymin=-.2,ymax=1.5,
                                xtick={-3,-2,-1,0,1,2,3}, ytick=\empty,
                                xticklabels={$-3f_s$,$-2f_s$,$-f_s$,$0$,$f_s$,$2f_s$,$3f_s$},
                                clip=false,
                            ]
                        \addplot+ [mark=none, solid,very thick, teal, domain=(-3.5:3.5), samples=100] {sinc(pi*x/2)^4} node[midway,above right] {\small $\frac{A^2}{16}T_s$};
                        \draw[very thick, -stealth, teal] (axis cs:0,0) -- (axis cs:0,1) node[midway,right] {\small $\frac{A^2}{16}$};
                        \draw[very thick, -stealth, teal] (axis cs:1,0) -- (axis cs:1,.16) node[midway,above right] {\small $\frac{A^2}{4\pi^2}$};
                        \draw[very thick, -stealth, teal] (axis cs:3,0) -- (axis cs:3,0.05);
                        \draw[very thick, -stealth, teal] (axis cs:-1,0) -- (axis cs:-1,.16) node[midway,above left] {\small $\frac{A^2}{4\pi^2}$};
                        \draw[very thick, -stealth, teal] (axis cs:-3,0) -- (axis cs:-3,0.05);
                        \end{axis}
                \end{tikzpicture}
        \end{center}

        \item Para el segundo caso, el ancho del pulso es $2$ tiempos de símbolo, $T=2T_s$, por lo que la onda PAM requerida queda:
        \begin{center}
                \begin{tikzpicture}
                        \begin{axis}[   x=1.8cm, y=1.5cm, axis lines=middle,  
                                xlabel=$t$,
                                ylabel=$x(t)$,
                                x label style={anchor=north},
                                y label style={anchor=east},
                                xmin=-1,xmax=7,ymin=-.2,ymax=1.4,
                                xtick={0,1,2,3,4,5,6}, ytick={1},
                                xticklabels={$0$,$T_s$,$2T_s$,$3T_s$,$4T_s$,$5T_s$,$6T_s$},
                                yticklabels={$A$},
                                clip=false,
                            ]
                        \addplot+[mark=none, dashed, red!50!white] coordinates {(-1,0) (0,1) (1,0)};
                        \addplot+[mark=none, dashed, red!50!white] coordinates {(0,0) (1,1) (2,0)};
                        \addplot+[mark=none, dashed, red!50!white] coordinates {(1,0) (2,1) (3,0)};
                        \addplot+[mark=none, dashed, red!50!white] coordinates {(2,0) (3,1) (4,0)};
                        \addplot+[mark=none, dashed, red!50!white] coordinates {(3,0) (4,1) (5,0)};
                        \addplot+[mark=none, dashed, red!50!white] coordinates {(4,0) (5,1) (6,0)};
                        \addplot+[mark=none, dashed, red!50!white] coordinates {(5,0) (6,1) (7,0)};

                        \addplot+[mark=none, solid, very thick, teal] coordinates {(-1,0) (0,1) (1,0) (2,1) (5,1) (6,0) (7,0)};
                        \end{axis}
                \end{tikzpicture}                        
        \end{center}
        Notar que de todos modos no hay ISI en el instante exacto de muestreo.

        Para calcular la densidad espectral de potencia, los valores $\mu$ y $\sigma^2$ son los anteriores. Además, $|\hat{P}(f)|^2 = T_s^2\sinc^4\left(\pi T_s f\right)$. De esto se deduce que:
        \begin{itemize}
                \item El espectro tiene una parte continua, que se anula en los múltiplos de $f_s$.
                \item La parte discreta incluye un valor de continua, pero no otras deltas ya que $|P(kf_s)|=0$ (similar al pulso cuadrado NRZ de las notas).
        \end{itemize}
        La expresión completa es:
        \begin{equation*}
                S_x(f) = \frac{A^2}{4}T_s \sinc^4\left(\pi\frac{f}{f_s}\right) + \frac{A^2}{4} \delta(0).
        \end{equation*}

        \begin{center}
                \begin{tikzpicture}
                        \pgfmathdeclarefunction{sinc}{1}{%
                        \pgfmathparse{(#1==0 ? 1: sin(#1 r))/(#1==0 ? 1: #1)}%
                        }
                        \begin{axis}[   x=1.8cm, y=1.5cm, axis lines=middle,  
                                xlabel=$f$,
                                ylabel=$S_x(f)$,
                                x label style={anchor=north},
                                y label style={anchor=east},
                                xmin=-3.5,xmax=3.5,ymin=-.2,ymax=1.5,
                                xtick={-3,-2,-1,0,1,2,3}, ytick=\empty,
                                xticklabels={$-3f_s$,$-2f_s$,$-f_s$,$0$,$f_s$,$2f_s$,$3f_s$},
                                clip=false,
                            ]
                        \addplot+ [mark=none, solid,very thick, teal, domain=(-3.5:3.5), samples=100] {sinc(pi*x)^4} node[midway,above right] {\small $\frac{A^2}{4}T_s$};
                        \draw[very thick, -stealth, teal] (axis cs:0,0) -- (axis cs:0,1) node[midway,right] {\small $\frac{A^2}{4}$};
                        \end{axis}
                \end{tikzpicture}
        \end{center}
        Observación, en ambos casos como se trata de $\sinc^4$ en frecuencia, los lóbulos menores llevan muy poca potencia.
\end{enumerate}

\end{solucion}

\begin{ejercicio}
        Sea la onda PAM $x(t)=\sum_k a_k p(t-k T_s)$ con pulsos $p(t)=\sinc(\pi f_s t)$. Los símbolos $a_k \in \{0, A, 2A, 3A\}$ son sorteados independientes y con probabilidades  respectivas 
        $\frac{1}{2}, \frac{1}{4}, \frac{1}{8}, \frac{1}{8}.$
        
        \begin{enumerate}
        \item Si el canal tiene transferencia $\hat{H}(f)$ como en la figura:

        \begin{center}
        \begin{tikzpicture}
          \draw [->,thick] (-4,0) -- (4,0);
          \draw [->,thick] (0,-0.5) -- (0,3);
          \draw[teal, thick] (-4,1) -- (-2,1) -- (-2,2) -- (2,2) -- (2,1) -- (4,1);
          \draw [dotted] (2,1) -- (2,0);
          \draw [dotted] (-2,1) -- (-2,0);
          \draw [dotted] (-2,1) -- (2,1);
          \draw (4.3,0) node {$f$};
          \draw (0,3.3) node {$\hat{H}(f)$};
          \draw (2,-0.3) node {$\frac{f_s}{3}$};
          \draw (-2,-0.3) node {$-\frac{f_s}{3}$};
          \draw (-0.2,1.3) node {$\frac{1}{2}$};
          \draw (-0.2,2.3) node {$1$};
        \end{tikzpicture}
        \end{center}

         Demostrar que en este caso se tiene ISI al detectar con la frecuencia $f_s$. Si ahora se transmite la señal $\tilde{x}(t)=\sum_k a_k p(t-k T'_s)$ con el mismo pulso $p(t)$, hallar $T'_s$ para que no haya ISI en la detección. A partir de ahora trabajamos con la señal $\tilde{x}(t)$.
        \item Hallar y bosquejar la densidad espectral de potencia de onda PAM recibida luego de pasar por el canal.
        \item ¿Es posible obtener sincronismo a partir de la señal recibida? Si las probabilidades fueran diferentes, ¿es posible que haya forma de sincronizar el detector? 
        \item Hallar la eficiencia espectral del sistema y compararla con el caso en que el canal fuera ideal ($\hat{H}(f) = 1 \,\, \forall f$) y pudiéramos usar la frecuencia $f_s$ original.
        \item Sabiendo de antemano que el canal tiene la transferencia de la primera parte y sin modificar la forma de los pulsos, ¿Es posible modificar el receptor de manera que no haya ISI utilizando la frecuencia $f_s$ original? Explicar.
        \end{enumerate}
        
\end{ejercicio}

\begin{solucion}

        \begin{enumerate}
        
              \item Comencemos analizando $p(t)$ en frecuencia:
                \begin{equation*}
                        \hat{P}(f) = \begin{cases}
                        T_s & -\frac{f_s}{2} \leqslant f \frac{f_s}{2} \\
                        0 & \textrm{fuera de la banda}
                        \end{cases}
                \end{equation*}
                Es decir un pulso de ancho de banda $\pm f_s/2$ en frecuencia. La onda PAM observada a la salida contiene al pulso $\tilde{p} = p*h$, deformado por el canal, o bien en frecuencia:
                \begin{center}
                        \begin{tikzpicture}
                                \begin{axis}[x=1.8cm, y=1.5cm, axis lines=middle,  
                                         xlabel=$f$,
                                         ylabel=$\hat{\tilde{P}}(f)$,
                                         x label style={anchor=north},
                                         y label style={anchor=east},
                                         xmin=-3.5,xmax=3.5,ymin=-.2,ymax=1.5,
                                         xtick={-3,-2,0,2,3}, ytick=\empty,
                                         xticklabels={$-f_s/2$,$-f_s/3$,$0$,$f_s/3$,$f_s/2$},
                                         clip=false,
                                     ]
                                \addplot+ [mark=none, solid,very thick, teal, domain=(-3.5:3.5), const plot] coordinates {
                                        (-3.5,0) (-3,.5) (-2,1) (0,1) (2,.5) (3,0) (3.5,0)} node[midway,above right] {$T_s$};
                                \node[teal, above] at (axis cs:2.5,.5) {$T_s/2$};
                                \end{axis}
                        \end{tikzpicture}
                \end{center}
                Observar que el pulso no tiene simetría vestigial alrededor de $f_s/2$, por lo que introducirá ISI al pasar por el canal si se utiliza tiempo de símbolo $T_s$. Sin embargo, sí posee simetría vestigial alrededor de $(5/12)f_s$, por lo que si tomamos $f'_s/2 = (5/12)f_s$, tendremos un pulso sin ISI. Esto corresponde a usar $f'_s = (5/6)f_s$.

                \item De las probabilidades dadas se tiene:
                \begin{align*}
                        \mu &= 0 (1/2) + A (1/4) + 2A (1/8) + 3A (1/8) = \frac{7}{8}A  \Rightarrow \mu^2 = \frac{49}{64}A^2\\
                        E[a_k^2] &= 0 (1/2) + A^2 (1/4) + 4A^2 (1/8) + 9A^2 (1/8) = \frac{15}{8}A^2\\
                        \sigma^2 &= \frac{15}{8}A^2 - \frac{49}{64}A^2 = \frac{71}{64}A^2
                \end{align*}
                El espectro $S_x(f)$ depende de $|\hat{\tilde{P}}|^2$ que es simplemente elevar al cuadrado las dos zonas horizontales. Notar que $f'_s T_s = f'_s/f_s = 5/6$. Esto modifica las alturas. Para la parte discreta del espectro, no aparecen armónicos ya que $(5/6)f_s > f_s/2$ por lo que $\hat{\tilde{P}}$ vale $0$. Solo aparece continua. El resultado final es:
                \begin{equation*}
                        S_x(f) = \sigma^2 f'_s |\hat{\tilde{P}}(f)|^2 + \mu^2 {f'_s}^2 T_s^2\delta(0)
                \end{equation*}

                \begin{center}
                        \begin{tikzpicture}
                                \begin{axis}[x=1.8cm, y=1.5cm, axis lines=middle,  
                                         xlabel=$f$,
                                         ylabel=$S_x(f)$,
                                         x label style={anchor=north},
                                         y label style={anchor=east},
                                         xmin=-3.5,xmax=3.5,ymin=-.2,ymax=1.7,
                                         xtick={-3,-2,0,2,3}, ytick=\empty,
                                         xticklabels={$-f_s/2$,$-f_s/3$,$0$,$f_s/3$,$f_s/2$},
                                         clip=false,
                                     ]
                                \addplot+ [mark=none, solid,very thick, teal, domain=(-3.5:3.5), const plot] coordinates {
                                        (-3.5,0) (-3,.25) (-2,1) (0,1) (2,.25) (3,0) (3.5,0)} node[midway,above right] {$0.92 A^2 T_s$};
                                \node[teal, above] at (axis cs:2.5,.25) {$0.23 A^2T_s$};
                                \draw[very thick, -stealth, teal] (axis cs:0,0) --(axis cs:0,.8) node[midway,right] {$0.53 A^2$};
                                \end{axis}
                        \end{tikzpicture}
                \end{center}

                \item No es posible obtener sincronismo, y como el espectro no llega hasta $f'_s$, no hay ninguna forma de obtenerlo cambiando las probabilidades de los símbolos.
                
                \item La $\textrm{vti}$ es $\log_2(M) f'_s = \frac{5}{3} f_s$. El ancho de banda del canal es $B=f_s/2$ por lo que la eficiencia espectral es $\frac{\textrm{vti}}{B} = \frac{(5/3)f_s}{f_s/2} = 3.33 \unit{\bit/\hertz} < 4 \unit{\bit/\hertz}$ que es el valor ideal en este caso.
                
                \item No a menos que se ecualice el canal en primera instancia.
        \end{enumerate}
\end{solucion}



\begin{ejercicio}
Sea la onda PAM $x(t)=\sum_k a_k p(t-k T_s)$ con el pulso
\begin{equation*}
        p(t) = 2B \left[ \sinc(\pi B t) \right]^2 \cos( \pi B t )       
\end{equation*}

\begin{enumerate}
\item Determinar el mínimo valor de $T_s$ en función de $B$ para que no haya ISI. A partir de ahora se utiliza ese valor.
\item Si $a_k$ puede tomar los valores $0$, $A$ y $2A$ con igual probabilidad e independencia entre un sorteo y otro, hallar y bosquejar la densidad espectral de potencia de onda PAM resultante. ¿Se puede obtener una señal para sincronizar el detector a partir de la onda recibida?
\item Hallar la potencia de la onda y la energía media del pulso transmitido.
\item Hallar la eficiencia espectral. Comparar con la que se hubiera obtenido usando pulsos de Nyquist ideales con el mismo ancho de banda.
\end{enumerate}
\end{ejercicio}

\begin{solucion}
        \begin{enumerate}
        \item Comencemos analizando el pulso en frecuencia. El término $B \sinc^2 (\pi B t)$ corresponde a un espectro triangular en frecuencia, de ancho total $2B$ y alto $1$ (en gris). El término $2\cos(\pi B t)$ corresponde a dos deltas de alto $1$ en $\pm B/2$ de frecuencia. Convolucionando el triángulo con cada delta (en rojo) y sumando queda:
        
        \begin{center}
                \begin{tikzpicture}
                        \begin{axis}[x=1.8cm, y=1.5cm, axis lines=middle,  
                                 xlabel=$f$,
                                 ylabel=$\hat{P}(f)$,
                                 x label style={anchor=north},
                                 y label style={anchor=east},
                                 xmin=-3.5,xmax=3.5,ymin=-.2,ymax=1.5,
                                 xtick={-3,-2,-1,0,1,2,3}, ytick=\empty,
                                 xticklabels={$-3B/2$, $-B$, $-B/2$,$0$,$B/2$,$B$, $3B/2$},
                                 clip=false,
                             ]
                        \addplot+ [mark=none, solid, gray] coordinates {
                                (-3.5,0) (-2,0) (0,1) (2,0) (3.5,0)
                        };
                        \draw[thick,gray,-stealth] (axis cs:-1,0) -- (axis cs:-1,1);
                        \draw[thick,gray,-stealth] (axis cs:1,0) -- (axis cs:1,1);
                        \addplot+ [mark=none, dashed, red!50!white] coordinates {
                                (-3.5,0) (-3,0) (-1,1) (1,0) (3.5,0)
                        };
                        \addplot+ [mark=none, dashed, red!50!white] coordinates {
                                (-3.5,0) (-1,0) (1,1) (3,0) (3.5,0)
                        };
                        \addplot+ [mark=none, very thick, solid, teal] coordinates {
                                (-3.5,0) (-3,0) (-1,1) (1,1) (3,0) (3.5,0)
                        } node[midway, above right] {$1$};
                        \end{axis}
                \end{tikzpicture}
        \end{center}
        Notemos que el pulso presenta simetría vestigial alrededor de $B$, por lo que podemos tomar $f_s = 2B$ para no tener ISI. Si queremos visualizarlo en el tiempo, notemos que $\sinc\left(\pi B \frac{k}{2B}\right)$ se anula en los múltiplos pares de $k$, mientras que $\cos\left(\pi B \frac{k}{2B}\right)$ se anula en los múltiplos impares. Entre ambos nos permiten usar todos los múltiplos de $T_s=\frac{1}{2B}$ sin ISI.

        \begin{center}
                \begin{tikzpicture}
                        \pgfmathdeclarefunction{sinc}{1}{%
                        \pgfmathparse{(#1==0 ? 1: sin(#1 r))/(#1==0 ? 1: #1)}%
                        }
                        \begin{axis}[   x=1.8cm, y=2cm, axis lines=middle,  
                                xlabel=$f$,
                                ylabel=$p(t)$,
                                x label style={anchor=north},
                                y label style={anchor=east},
                                xmin=-3.5,xmax=3.5,ymin=-.2,ymax=1.5,
                                xtick={-3,-2,-1,0,1,2,3}, ytick=\empty,
                                xticklabels={$-\frac{3}{2B}$,$-\frac{1}{B}$, $-\frac{1}{2B}$,$0$,$\frac{1}{2B}$,$\frac{1}{B}$,$\frac{3}{2B}$},
                                clip=false,
                            ]
                        \addplot+ [mark=none, solid,very thick, teal, domain=(-3.5:3.5), samples=100] {sinc(pi*x/2)^2*cos(deg(pi*x/2))} node[midway,above right] {$2B$};
                        \end{axis}
                \end{tikzpicture}
        \end{center}

        \item Calculamos primero:
         \begin{align*}
                \mu &= 0 (1/3) + A (1/3) + 2A (1/3) = A \Rightarrow \mu^2 = A^2\\
                E[a_k^2] &= 0 (1/3) + A^2 (1/3) + 4A^2 (1/3) = (5/3)A^2 \\
                \sigma^2 &= (2/3)A^2
         \end{align*}

         El espectro tendrá una parte continua, correspondiente a elevar al cuadrado el trapecio (ver ej. 1.5 para un caso similar). Como además, $\mu\neq 0$, habrá un término de continua. Como el espectro del pulso finaliza antes de $f_s = 2B$, no habrá otros armónicos.

         Bosquejo:
         \begin{center}
                \begin{tikzpicture}
                        \begin{axis}[x=1.8cm, y=1.5cm, axis lines=middle,  
                                xlabel=$f$,
                                ylabel=$S_x(f)$,
                                x label style={anchor=north},
                                y label style={anchor=east},
                                xmin=-3.5,xmax=3.5,ymin=-.2,ymax=1.5,
                                xtick={-3,-2,-1,0,1,2,3}, ytick=\empty,
                                xticklabels={$-3B/2$, $-B$, $-B/2$,$0$,$B/2$,$B$, $3B/2$},
                                clip=false,
                            ]
                        \addplot+[very thick,mark=none, const plot, teal, domain=(-1:1)] {1} node[above right,midway]{$A^2/(3B)$};
                        \addplot+[very thick,mark=none, teal, domain=(-3:-1)] {(x+3)^2/4};
                        \addplot+[very thick,mark=none, teal, domain=(1:3)] {(3-x)^2/4};
                        \draw[very thick, -stealth, teal] (axis cs:0,0) --(axis cs:0,.8) node[midway,right] {$4A^2B^2$};
                        \end{axis}
                \end{tikzpicture}              
        \end{center}

        Se tiene que:
        \begin{equation*}
                S = \int_{-\infty}^\infty S_x(f)df = \left(\frac{1}{3}\right)\frac{A^2}{3B}B + \frac{A^2}{3B}B+ \left(\frac{1}{3}\right)\frac{A^2}{3B}B + 4A^2B^2= \left(\frac{5}{9}\right) A^2 + 4A^2 B^2.
        \end{equation*}

        Para la energía media del pulso transmitido, necesitamos calcular $\bar{E}_p = E[a_k^2]||p||^2$. Para calcular $||p||^2$ conviene hacerlo en frecuencia elevando al cuadrado $|\hat{P}(f)|$ obteniendo:
        \begin{equation*}
                \bar{E}_p = (5/3)A^2 [(1/3)B + B + (1/3)B] = \frac{25}{9}A^2B
        \end{equation*}

        \item $\mathrm{vti} = f_s \log_2(M) = 2B \log_2(3)$. Ancho de banda $3B/2$, de donde la eficiencia es:
        \begin{equation*}
                \textrm{Eff} = \frac{2B\log_2(3)}{(3/2)B} = \frac{4}{3}\log_2(3) \approx 2.11 \unit{\bit/\hertz} < 3.17 \approx 2\log_2(3),
        \end{equation*}
        que es la eficiencia del pulso ideal en este caso.
        \end{enumerate}
\end{solucion}

\begin{ejercicio}
Sea la onda PAM $x(t)=\sum_k a_k p(t-k T_s)$, con el pulso
\begin{equation*}
        p(t) = 3B  \sinc(\pi B t)  \sinc(\alpha\pi B t).        
\end{equation*}
Para todas las partes salvo la 5., se utilizará $\alpha=\frac{1}{2}$.

\begin{enumerate}
\item Hallar $P(f)$, y determinar el mínimo valor de $T_s$ en función de $B$ para que no haya ISI. En lo que sigue se utilizara ese valor.
\item Si $a_k$ puede tomar los valores $-A,A$ en forma independiente y equiprobable, hallar y bosquejar la densidad espectral de potencia de onda PAM resultante. ¿Se puede obtener una señal para sincronizar el detector a partir de la onda recibida? Justificar.
\item Si ahora $a_k$ toma valores independientes $\{-A,A\}$ pero con probabilidades $1/3$ y $2/3$ respectivamente, cambia la respuesta a la parte 2? Justificar.
\item Hallar la potencia de la onda, y la energía media del pulso transmitido para las señales de 2 y 3.
\item Se considera ahora $\alpha>0$ cualquiera. Hallar $P(f)$ y determinar el mínimo valor de $T_s$ en función de $B$ para que no haya ISI. Discutir el tema de la sincronización del detector en este caso.
\end{enumerate}

\end{ejercicio}

\begin{solucion}
Graficamos primero el pulso, para $\alpha=1/2$:
\begin{center}
        \begin{tikzpicture}
                \pgfmathdeclarefunction{sinc}{1}{%
                \pgfmathparse{(#1==0 ? 1: sin(#1 r))/(#1==0 ? 1: #1)}%
                }
                \begin{axis}[   x=1.8cm, y=1cm, axis lines=middle,  
                        xlabel=$f$,
                        ylabel=$p(t)$,
                        x label style={anchor=north},
                        y label style={anchor=east},
                        xmin=-3.5,xmax=3.5,ymin=-.2,ymax=3.5,
                        xtick={-3,-2,-1,0,1,2,3}, ytick=\empty,
                        xticklabels={$-\frac{3}{B}$,$-\frac{2}{B}$, $-\frac{1}{B}$,$0$,$\frac{1}{B}$,$\frac{2}{B}$,$\frac{3}{B}$},
                        clip=false,
                    ]
                \addplot+ [mark=none, solid,very thick, teal, domain=(-3.5:3.5), samples=100] {3*sinc(pi*x)*sinc(pi*x/2)} node[midway,above right] {$3B$};
                \end{axis}
        \end{tikzpicture}
\end{center}

\begin{enumerate}
\item Para hallar $\hat{P}(f)$, conviene escribir $p(t) = \frac{3}{\alpha B} [B\sinc(\pi B t)][\alpha B \sinc(\alpha\pi B t)]$. De este modo, en frecuencia es la convolución de dos pulsos de anchos $B$ y $\alpha B$ (similar al Ej. 1.5) y alto $1$. Esto da el siguiente trapecio (para $\alpha=1/2$):
\begin{center}
        \begin{tikzpicture}
                \begin{axis}[   x=1.8cm, y=1.5cm, axis lines=middle,  
                        xlabel=$f$, ylabel=$\hat{P}(f)$,
                        x label style={anchor=north},
                        y label style={anchor=east},
                        xmin=-3.5,xmax=3.5,ymin=-.2,ymax=1.4,
                        xtick={-3,-2,-1,0,1,2,3}, ytick=\empty,
                        xticklabels={$-\frac{3}{4}B$, $-\frac{1}{2}B$, $-\frac{1}{4}B$, $0$, $\frac{1}{4}B$, $\frac{1}{2}B$, $\frac{3}{4}B$},
                        clip=false,
                    ]
                \addplot[very thick,mark=none, teal] coordinates
                {(-3.5,0) (-3,0) (-1,1) (0,1) (1,1) (3,0) (3.5,0)} node[above right,midway]{$3$};
                \end{axis}
        \end{tikzpicture}                        
\end{center}
Notar que el área bajo el trapecio es $3B = p(0)$ como verificación. El pulso en frecuencia tiene simetría vestigial alrededor de $(1/2)B$ por lo que podemos tomar $f_s/2 = (1/2)B$ o bien $f_s=B$.

\item En este caso $\mu=0$ y $\sigma^2 = E[a_k^2]= A^2$. La onda recibida tendrá solo densidad continua en frecuencia, sin deltas.

Bosquejo:
         \begin{center}
                \begin{tikzpicture}
                        \begin{axis}[x=1.8cm, y=1.5cm, axis lines=middle,  
                                xlabel=$f$,
                                ylabel=$S_x(f)$,
                                x label style={anchor=north},
                                y label style={anchor=east},
                                xmin=-3.5,xmax=3.5,ymin=-.2,ymax=1.5,
                                xtick={-3,-2,-1,0,1,2,3}, ytick=\empty,
                                xticklabels={$-3B/4$, $-B/2$, $-B/4$,$0$,$B/4$,$B/2$, $3B/4$},
                                clip=false,
                            ]
                        \addplot+[very thick,mark=none, const plot, teal, domain=(-1:1)] {1} node[above right,midway]{$9A^2 B$};
                        \addplot+[very thick,mark=none, teal, domain=(-3:-1)] {(x+3)^2/4};
                        \addplot+[very thick,mark=none, teal, domain=(1:3)] {(3-x)^2/4};
                        \end{axis}
                \end{tikzpicture}              
        \end{center}


\item Si cambian las probabilidades, aparecerá una $\delta$ en continua, de valor $\mu^2 B^2$, pero no armónicos ya que el espectro del pulso termina antes de $f_s=B$. En este caso $\mu^2 = A^2/9$.

\item Para el caso $\mu=0$, hay que integrar la $S_x(f)$ resultando en:
 \begin{equation*}
        S = 9A^2 B \left[\frac{B}{6} + \frac{B}{2} + \frac{B}{6} \right] = \frac{15}{2}A^2 B^2
 \end{equation*}

 Si además hay continua, hay que sumar la $\delta$, obteniendo $S=\frac{15}{2}A^2 B^2 + \frac{1}{9}A^2B^2$.

 La energía media del pulso transmitido surge de $\bar{E}_p = E[a_k^2]||p||^2$. Integrando en frecuencia, $||p||^2 = 9\left[\frac{B}{6} + \frac{B}{2} + \frac{B}{6} \right] = \frac{15}{2}B$, de donde $\bar{E}_p = \frac{15}{2}A^2B$. Como $E[a_k^2]$ no depende de las probabilidades, la energía media del pulso es la misma en ambos casos.

 \item Si ahora $\alpha>0$ cualquiera, tendremos un trapecio como el anterior si $0<\alpha<1$. Para $\alpha=1$ queda un triángulo de ancho $2B$ y podemos seguir usando $f_s=B$ para tener simetría vestigial. Si ahora $\alpha>1$, el pulso ``lento'' es el que tiene ancho $B$ (ocupa menos ancho de banda) y el espectro vuelve a ser un trapecio de altura $3/\alpha$ ahora, y con simetría vestigial en $f_s/2 = \alpha B/2$, es decir que podemos aumentar la frecuencia de señalización a $f_s=\alpha B$ (correspondiente a las raíces del segundo $\sinc$) a cambio de mayor ancho de banda. De todos modos, el espectro siempre está contenido dentro de $\pm f_s$ por lo que en caso de haber continua no habrá armónicos.
\end{enumerate}

\end{solucion}

\section{Error en detección sincrónica en presencia de ruido}

\begin{ejercicio}
Se considera un sistema de trasmisión digital con señalización binaria unipolar $a_k=\{0,A\}$, pulsos rectangulares NRZ, y donde el ruido en el detector es Gaussiano con media nula y varianza $\sigma^2$. Si la probabilidad de error de símbolo es $P_e=1\times 10^{-5}$, calcular la relación $\left(\frac{A}{\sigma}\right)$. Si en cambio se usa señalización polar con igual relación señal/ruido, ¿cuál es la probabilidad de error de símbolo?        
\end{ejercicio}

\begin{solucion}
Si se usa señalización polar, el umbral es $A/2$ y necesitamos $Q(A/(2\sigma)) = 10^{-5}$. De la gráfica, $A/(2\sigma) = 4.26$ o bien $A/\sigma = 8.52$. Para señalización polar, con este $A/\sigma$, el umbral es $0$ y la prob. de error es $Q(A/\sigma) = Q(8.52) \approx 7.3\times 10^{-18}$.
\end{solucion}


\begin{ejercicio}
Se considera un sistema de transmisión con señalización binaria polar, donde existe interferencia intersimbólica (ISI). Modelamos esta interferencia como una tensión $\pm \alpha A$ que se añade al valor de señal en el instante de muestreo, con $0 < \alpha < 1$. Suponiendo que tanto los símbolos $\pm A$ como los valores de interferencia $\pm \alpha A$ son equiprobables e independientes entre sí, determinar la probabilidad de error de símbolo. En el caso $A=8\sigma$ y $\alpha=0.1$ comparar las probabilidades de error con y sin ISI.
\end{ejercicio}

\begin{solucion}
        Notemos que los símbolos posibles a recibir van a ser $\{-(1+\alpha)A, -(1-\alpha)A, (1-\alpha)A, (1+\alpha)A\}$, todos con igual probabilidad (todas las combinaciones de símbolo e ISI modelada). Al detectar tendremos un ruido Gaussiano de potencia $\sigma^2$. Colocamos el umbral el $0$ por simetría y la probabilidad de error queda:
        \begin{equation*}
                P_e = \frac{Q((1-\alpha)(A/\sigma))+Q((1+\alpha)(A/\sigma))}{2}
        \end{equation*}

        Para $A/\sigma=8$ y $\alpha=0.1$ queda:
        \begin{equation*}
                P_e = \frac{Q(7.2)+Q(8.8)}{2} \approx 1.5\times 10^{-13}.
        \end{equation*}
        Para $\alpha=0$ (sin ISI), queda $Q(8)= 6\times 10^{-16}$. La ISI aumenta mucho la probabilidad de error, de hecho, la mitad del tiempo estamos trabajando con un un $A$ efectivo de $(1-\alpha)A$ y en general $Q((1+\alpha)(A/\sigma)) \ll Q((1-\alpha)(A/\sigma))$ por lo que $P_e \approx Q((1-\alpha)(A/\sigma))/2$.
\end{solucion}

\begin{ejercicio}
Un modelo comúnmente utilizado para el ruido impulsivo en los canales telefónicos, es ruido blanco con función de densidad de Laplace o doble exponencial

\bigskip

\begin{minipage}{0.4\textwidth}
\begin{equation*}
f_N(n)= \frac{1}{\sqrt{2\sigma^2}} e^{-\sqrt{2}|n|/\sigma}
\end{equation*}
\end{minipage}%
\begin{minipage}{0.6\textwidth}
\begin{center}
\begin{tikzpicture}
        \begin{axis}[   x=1cm, y=1.5cm, axis lines=middle,  
                xlabel=$n$,
                x label style={anchor=north west},
                y label style={anchor=east},
                xmin=-3.2,xmax=3.2,ymin=-.2,ymax=1.2,
                xtick=\empty, ytick=\empty,
                xticklabels={},
                clip=false,
                smooth
            ]
        \addplot+[teal, thick,mark=none,domain=0:3] {exp(-x)};
        \addplot+[teal, thick,mark=none,domain=-3:0] {exp(x)};
        \end{axis}
\end{tikzpicture}
\end{center}
\end{minipage}


\begin{enumerate}
\item Calcule la probabilidad de error de símbolo en un sistema de trasmisión digital binario que usa pulsos de Nyquist de amplitud $\pm A$, y donde el canal presenta este ruido impulsivo de varianza $\sigma^2$.
\item Si se desea obtener una probabilidad de error menor a $1\times 10^{-3}$, compare las amplitudes de los pulsos de Nyquist necesarias en cada caso (canal con ruido impulsivo y canal con ruido gaussiano).
\end{enumerate}
\end{ejercicio}

\begin{solucion}
\begin{enumerate}
\item El usar pulsos de Nyquist aquí no es más para indicar que son de ancho de banda acotado y no hay ISI, por lo que el único error se debe al ruido. Si transmito $A$, al perturbar por el ruido la densidad del valor recibido es:
\item 
        \begin{minipage}{0.4\textwidth}
        \begin{equation*}
        f_x(x\mid A)= \frac{1}{\sqrt{2\sigma^2}} e^{-\sqrt{2}|x-A|/\sigma}
        \end{equation*}
        \end{minipage}%
        \begin{minipage}{0.6\textwidth}
        \begin{center}
        \begin{tikzpicture}
                \begin{axis}[   x=1cm, y=1.5cm, axis lines=middle,  
                        xlabel=$n$,
                        x label style={anchor=north west},
                        y label style={anchor=east},
                        xmin=-3.2,xmax=3.2,ymin=-.2,ymax=1.2,
                        xtick={1}, ytick=\empty,
                        xticklabels={$A$},
                        clip=false,
                        smooth
                    ]
                \addplot+[teal, thick,mark=none,domain=1:3] {exp(-(x-1))};
                \addplot+[teal, thick,mark=none,domain=-3:1] {exp(x-1)};
                \end{axis}
        \end{tikzpicture}
        \end{center}
        \end{minipage}
Integrando:
\begin{equation*}
        P(x<0\mid A) = \int_{-\infty}^0 \frac{1}{\sqrt{2\sigma^2}} e^{-\sqrt{2}|x-A|/\sigma} dx = \frac{1}{\sqrt{2\sigma^2}} \int_{-\infty}^0 e^{-\frac{\sqrt{2}}{\sigma} (A-x)}dx = \frac{1}{2} e^{-\sqrt{2}\frac{A}{\sigma}}.
\end{equation*} 
Por simetría $P(x>0\mid -A)$ es igual, y por lo tanto la probabilidad de error es:
\begin{equation*}
        P_e = \frac{1}{2}P(x<0\mid A)+ \frac{1}{2}P(x>0\mid -A) = \frac{1}{2} e^{-\sqrt{2}\frac{A}{\sigma}}.
\end{equation*}

\item Para ruido impulsivo igualamos:
\begin{equation*}
        \frac{1}{2} e^{-\sqrt{2}\frac{A}{\sigma}} = 10^{-3} \Rightarrow \frac{A}{\sigma} \approx 4.4.
\end{equation*}
Comparando, para ruido Gaussiano necesitamos $Q(A/\sigma) = 10^{-3}$ de donde $A/\sigma = 3.1$.

\end{enumerate}

\end{solucion}


\begin{ejercicio}
Se recibe una señal $\sum_k a_k p(t-kT_s)$, con símbolos unipolares binarios $a_k
\in \{0,A\}$, y pulsos triangulares $ p(t)=1 - \frac{|t|}{T_s}$, contaminada con ruido blanco de densidad espectral $\frac{\eta_0}{2}$.
\begin{enumerate}
\item Especificar el detector óptimo para esta forma de pulso.
\item Calcular la probabilidad de error de símbolo en función de los parámetros $A$, $T_s$ y $\eta$.
\end{enumerate} 
\end{ejercicio}

\begin{solucion}
        Al ser un pulso triangular, el filtro acoplado correspondiente (causal) es otro triángulo $h(t) = 1-\frac{|t-T_s|}{T_s}$ entre $0$ y $2T_s$, y muestrear con $\tau=T_s$ de retardo. Podemos además calcular $||p||^2$ integrando en el tiempo para obtener $||p||^2 = \frac{2}{3}T_s$ y ajustar la ganancia $\kappa=1/||p||^2 = \frac{3}{2T_s}$ para recuperar el símbolo exacto a la salida.
        \begin{center}
                \begin{tikzpicture}
                        \begin{axis}[   x=3cm, y=1.5cm, axis lines=middle,  
                                xlabel=$t$,
                                ylabel=$p(t)$,
                                x label style={anchor=north},
                                y label style={anchor=east},
                                xmin=-1.2,xmax=1.2,ymin=-.2,ymax=1.4,
                                xtick={-.3,0,.3}, ytick={1},
                                xticklabels={$-T_s$, 0 , $T_s$},
                                clip=false,
                            ]
                        \addplot+[very thick,mark=none, teal] coordinates {(-1,0) (-.3,0) (0,1) (.3,0) (1,0)};
                        \end{axis}
                \end{tikzpicture} \hspace{1em}
                \begin{tikzpicture}
                        \begin{axis}[   x=3cm, y=1.5cm, axis lines=middle,  
                                xlabel=$t$,
                                ylabel=$h(t)$,
                                x label style={anchor=north},
                                y label style={anchor=east},
                                xmin=-1.2,xmax=1.2,ymin=-.2,ymax=1.4,
                                xtick={0,.3,.6}, ytick={1},
                                xticklabels={0, $T_s$, $2T_s$}, yticklabels = {$\frac{3}{2T_s}$},
                                clip=false,
                            ]
                        \addplot+[very thick,mark=none, teal] coordinates {(-1,0) (0,0) (0.3,1) (.6,0) (1,0)};
                        \end{axis}
                \end{tikzpicture}      
        \end{center}

        Por ser filtro acoplado, la relación señal a ruido a la salida es:
        \begin{equation*}
                \left(\frac{A}{\sigma}\right)^2_{\max} = 2 \frac{A^2}{\eta} ||p||^2 = \frac{4A^2T_s}{3\eta}.
        \end{equation*}

        Por ser unipolar binario, la probabilidad de error es entonces:
        \begin{equation*}
                P_e = Q\left(\frac{A}{2\sigma}\right) = Q\left(\sqrt{\frac{A^2 T_s}{3\eta}}\right)
        \end{equation*}

        Observación: el espectro del pulso es infinito en frecuencia, por lo que se requiere que el canal tenga suficiente ancho de banda como para no incorporar ISI (al menos $\pm f_s$ para dejar pasar el lóbulo central).
\end{solucion}

\begin{ejercicio}
Un sistema de comunicación digital binario utiliza un canal con ruido blanco gaussiano de densidad espectral de potencia $\eta/2=10^{-15}$ $V^2/\unit{\Hz}$. La velocidad de señalización es $1\times 10^6$ Baudios, y los pulsos son rectangulares NRZ.
\begin{enumerate}
\item Calcular la potencia necesaria para que la probabilidad de error de símbolo $P_e$ sea menor que $1\times 10^{-6}$ si se utiliza señalización unipolar. Recalcular para señalización polar.
\item ¿Cuál sería la potencia necesaria para mantener la misma velocidad de transmisión de información con igual probabilidad de error $P_e$ usando señalización M-aria, con una constelación de 8 símbolos de amplitudes $\{\pm 1,\pm 3,\pm 5,\pm 7\}$.
\end{enumerate}
\end{ejercicio}

\begin{solucion}
\begin{enumerate}
        \item Para el pulso NRZ, $||p||^2 = T_s = 1/f_s$ en este caso. De aquí se deduce que, usando filtro acoplado:
        \begin{equation*}
                \left(\frac{A}{\sigma}\right)^2_{\max} = 2 \frac{A^2}{\eta} ||p||^2 = \frac{2A^2}{f_s \eta},
        \end{equation*}
        y la probabilidad de error para el caso unipolar queda:
        \begin{equation*}
                P_e = Q\left(\sqrt{\frac{A^2}{2f_s \eta}}\right) = 10^{-6}.
        \end{equation*}
        De donde:
        \begin{equation*}
        \frac{A}{\sqrt{2f_s \eta}} \approx 4.75 \Rightarrow A \approx 213\mu V.
        \end{equation*}

        La energía media del pulso NRZ es $\bar{E}_p = \frac{A^2}{2}T_s = \frac{A^2}{2f_s}$ y la potencia total es $S = \bar{E}_p f_s = \frac{A^2}{2} \approx 2.26 \times 10^{-8} \approx -46 \mathrm{dBm}$.

        Si la codificación es polar binaria, la probabilidad de error es:
        \begin{equation*}
                P_e = Q\left(\sqrt{\frac{2A^2}{f_s \eta}}\right) = 10^{-6}.
        \end{equation*}
        De donde:
        \begin{equation*}
        A\sqrt{\frac{2}{f_s \eta}} \approx 4.75 \Rightarrow A \approx 106 \mu V.
        \end{equation*}

        La energía media del pulso NRZ es $\bar{E}_p = A^2 T_s = \frac{A^2}{f_s}$ y la potencia total es $S = \bar{E}_p f_s = A^2 \approx 1.13 \times 10^{-8} \approx -49 \mathrm{dBm}$ (la mitad).

        \item $\mathrm{vti} = 10^6 \unit{\bit/\sec}$, de donde podemos bajar $f_s$ a $f_s = \frac{\mathrm{vti}}{\log_2 M} \approx 333 \unit{\kilo\hertz}$. La probabilidad de error de bit pedida es $10^{-6}$ por lo que la probabilidad de error de símbolo es un poco más alta, $P_e(\mathrm{sym}) = \log_2(M) P_e(\mathrm{bit}) = 3\times 10^{-6}$.
        
        Al ser un sistema $M-$ario polar, la probabilidad de error de símbolo es:
        \begin{equation*}
                P_e(\mathrm{sym}) = 2\frac{M-1}{M}Q\left(\frac{A}{\sigma}\right) = 2\frac{M-1}{M}Q\left(\frac{A}{\sigma}\right) = \frac{14}{8}Q\left(\frac{A}{\sigma}\right) = 3\times 10^{-6}
        \end{equation*}
        De donde se requiere un $A/\sigma \approx 4.64$. Al ser filtro acoplado, el $(A/\sigma)^2$ máximo es $2\frac{A^2}{\eta}||p||^2= 2\frac{A^2}{\eta f_s}$ para el pulso NRZ. Por lo tanto:
        \begin{equation*}
        4.64^2 = 2\frac{A^2}{\eta f_s} \Rightarrow A \approx 60 \unit{\mu\volt}
        \end{equation*}

        La potencia transmitida es $\bar{E}_p f_s$, que en el caso $M-$ario polar queda:
        \begin{equation*}
                S = A^2 \frac{M^2-1}{3} \approx 7.55\times 10^{-8} \approx -41\mathrm{dBm}.
        \end{equation*}
        Requiere algo más de potencia en menor ancho de banda.
\end{enumerate}        
\end{solucion}


\begin{ejercicio}
Un sistema de transmisión con señalización binaria polar, pulsos de Nyquist ideales, y 50 repetidores debe tener probabilidad de error menor que $10^{-4}$. Calcular la relación $(S/N)$ a la entrada de cada repetidor cuando éstos son regenerativos. Repetir para el caso en que los repetidores sean no-regenerativos.
\end{ejercicio}

\begin{solucion}
        Tenemos $m=50$ repetidores regenerativos. La probabilidad de error en $m$ etapas debe ser $mp = 10^{-4}$ por lo que la probabilidad de error de etapa debe ser $p=2\times 10^{-6}$. De aquí $\sqrt{(S/N)_1}=4.61$ o bien $(S/N)_1=21.16 \approx 13$ dB.

        Si en cambio los repetidores son no regenerativos, debe ser $Q(\sqrt{(S/N)_m}) = 10^-4$ de donde $(S/N)_m = 13.83$ y por lo tanto $(S/N)_1= m(S/N)_m \approx 692 \approx 28$ dB.  
\end{solucion}

\begin{ejercicio}
Sea la onda digital $x_T(t)= \sum_k a_k p_N(t-kT_S)$, con símbolos $a_k \in \{ 0, A, 2A\}$ equiprobables, y pulsos de Nyquist $p_N$ de coseno alzado, y exceso de ancho de banda $\alpha=0.8$.

\begin{enumerate}
\item Calcular la eficiencia espectral de la comunicación en $\unit{bps/\hertz}$.
\item Si se recibe esta señal contaminada de ruido blanco, diseñar en diagrama de bloques un detector apropiado indicando todos los detalles relevantes. 
\item Para el detector diseñado, dar una fórmula para la probabilidad de error $P_e$ de símbolo en función de $A$, $f_s$ y la densidad espectral del ruido en el canal $\eta_0$. Evaluar $P_e$ para el caso de densidad espectral de ruido $\eta_0/2=10^{-15} \unit{\square\volt/\hertz}$, amplitud $A=0,6\,\unit{\milli\volt}$, y una velocidad de transmisión de información de 4 \unit{\mega\bit/\sec}.
\end{enumerate}
\end{ejercicio}

\begin{solucion}
\begin{enumerate}
        \item El pulso $p_N$ requiere un ancho de banda $B = \frac{f_s}{2}(1+\alpha) = 0.9f_s$. La $\mathrm{vti} = \log_2(3)f_s$ de donde la eficiencia queda:
        \begin{equation*}
                \mathrm{Eff} = \frac{\mathrm{vti}}{B} = \frac{\log_2(3)}{0.9} \approx 1.76.
        \end{equation*}

        \item Se debe usar un filtro acoplado, con un pulso de Nyquist de coseno alzado igual a $p_N$. Como este pulso es infinito en el tiempo, se debe tomar un $\tau$ grande como para incluir buena parte de la respuesta del filtro del lado causal (ej: $\tau>4T_s$) y lograr el máximo de acople. Se debe tener cuidado porque este filtro introduce ISI (al ser coseno alzado y caer como $t^3$, la ISI es pequeña y además no se requiere un $\tau$ tan alto).
        
        Si se usa el pulso $h(t) = p_N(\tau-t)/||p_N||^2$, en la recepción tendremos aproximadamente los símbolos $0,A,2A$ transmitidos. Conviene usar como umbrales $A/2$ y $3A/2$ para distinguir los símbolos.

        \item Para el sistema de detección anterior, la probabilidad de error de símbolo es:
        \begin{equation*}
                P_e(\mathrm{sym}) = \frac{1}{3}Q\left(\frac{A}{2\sigma}\right) + \frac{1}{3}2Q\left(\frac{A}{2\sigma}\right) + \frac{1}{3} Q\left(\frac{A}{2\sigma}\right) = \frac{4}{3}Q\left(\frac{A}{2\sigma}\right)
        \end{equation*}

        Al usar filtro acoplado tenemos que en recepción:
        \begin{equation*}
                \left(\frac{A}{\sigma}\right)^2_{\max} = 2 \frac{A^2}{\eta} ||p||^2
        \end{equation*}

        De donde la probabilidad de error queda:
        \begin{equation*}
                P_e(\mathrm{sym}) =  \frac{4}{3} Q\left(\sqrt{\frac{A^2}{\eta}||p||^2}\right)
        \end{equation*}

        Para completar la expresión, falta calcular la energía del pulso de Nyquist de coseno alzado con exceso $\alpha=0.8$, integrando en frecuencia $\int_{-\infty}^{\infty} |\hat{P}_N(f)|^2df$. Esto da:
        \begin{equation*}
                ||p||^2 = T_s\left(1-\frac{\alpha}{4}\right).
        \end{equation*}
        Sustituyendo:
        \begin{equation*}
                P_e(\mathrm{sym}) =  \frac{4}{3}Q\left(\sqrt{\frac{0.8 A^2}{\eta f_s}}\right)
        \end{equation*}

        Para los valores dados, se requiere $\textrm{vti} = 4\times 10^6$ por lo que $f_s = \textrm{vti}/\log_2(3) \approx 2.52 \unit{\mega\hertz}$. Sustituyendo queda:
        \begin{equation*}
                \sqrt{\frac{0.8 A^2}{\eta f_s}} = 7.55 \Rightarrow P_e = 4Q(7.55) = 2.9 \times 10^{-14}.
        \end{equation*}       

\end{enumerate}
\end{solucion}

\begin{ejercicio}
Se considera una onda PAM de la forma $\sum_k a_kp(t-kT_s)$, con símbolos $a_k\in\{A_1,A_2\}$, $A_1<A_2$, cuyas probabilidades, $P(A_1)$ y $P(A_2)$, no son necesariamente iguales.
        
\begin{enumerate}
\item Demuestre que el umbral óptimo para de detección para ruido Gaussiano es
\begin{equation*}
A^{\star} = \frac{A_1+A_2}{2} - \frac{\sigma^2\log\left(\frac{P(A_2)}{P(A_1)}\right)}{A_2-A_1},
\end{equation*}
y encuentre una expresión para la probabilidad de error $P_e$.

\item Supongamos que los pulsos $p(t)$ son pulsos de Nyquist de coseno alzado con exceso de ancho de banda 100\%. La señalización es polar con $A=500 \unit{\micro\volt}$ y velocidad $10 \unit{\mega\bit/\sec}$. Además, el canal presenta ruido blanco Gaussiano aditivo con densidad espectral de potencia $\eta_0/2~=~5\times 10^{-16} \,\unit{\volt}^2/\unit{Hz}$.

\begin{enumerate}
\item Si para la detección digital se utiliza un filtro acoplado, calcule la varianza del ruido a la salida del filtro. ¿Cuál es la desventaja de utilizar este tipo de filtro?

\item Proponga un filtro de recepción alternativo que solucione el inconveniente anterior. Utilizando el umbral óptimo de la parte 1, calcule la probabilidad de error $P_e$, para el caso $P(A_1)=\frac{1}{3}$, $P(A_2)=\frac{2}{3}$,

\end{enumerate}
\end{enumerate}
\end{ejercicio}

\begin{solucion}
\begin{enumerate}
\item La señal contaminada con ruido tiene la siguiente densidad, mezcla de dos Gaussianas:
        \begin{equation*}
                f_X(x) = P(A_1) \frac{1}{\sqrt{2\pi}\sigma}e^{-\frac{(x-A_1)^2}{2\sigma^2}} + P(A_2)\frac{1}{\sqrt{2\pi}\sigma}e^{-\frac{(x-A_2)^2}{2\sigma^2}}.
        \end{equation*}
        \begin{center}
        \begin{tikzpicture}
                \pgfmathdeclarefunction{gauss}{3}{%
                \pgfmathparse{1/(#3*sqrt(2*pi))*exp(-((#1-#2)^2)/(2*#3^2))}%
                }      
                \begin{axis}[x=2cm, y=5cm, axis lines=middle,  
                        xlabel=$t$,
                        ylabel=$f_X(x)$,
                        x label style={anchor=north},
                        y label style={anchor=east},
                        xmin=-1,xmax=4,ymin=-.2,ymax=.8,
                        xtick={1,1.894,3}, ytick=\empty,
                        xticklabels={$A_1$, $A^*$, $A_2$},
                        clip=false,
                    ]
                \addplot+[dashed ,mark=none, red, domain=(-1:4), samples=100] {.3*gauss(x,1,.5)};
                \addplot+[dashed ,mark=none, red, domain=(-1:4), samples=100] {.7*gauss(x,3,.5)};
                \addplot+[very thick,mark=none, teal, domain=(-1:4), samples=100] {.3*gauss(x,1,.5)+.7*gauss(x,3,.5)};
                \end{axis}
        \end{tikzpicture}
        \end{center}
        El umbral óptimo es cuando se igualan las densidades de cada símbolo, es decir:
        \begin{equation*}
                \frac{P(A_2)}{P(A_1)} = \frac{e^{-\frac{(A^*-A_1)^2}{2\sigma^2}}}{e^{-\frac{(A^*-A_2)^2}{2\sigma^2}}} = e^{\frac{(A^*-A_2)^2-(A^*-A_1)^2}{2\sigma^2}}.
        \end{equation*}
        Por lo tanto:
        \begin{equation*}
                2\sigma^2 \log\left(\frac{P(A_2)}{P(A_1)}\right) = (A^*-A_2)^2-(A^*-A_1)^2 = (2A^*-A_1-A_2)(A_1-A_2),
        \end{equation*}
        de donde despejando $A^*$ se obtiene la fórmula pedida.

        La probabilidad de error es:
        \begin{equation*}
                P_e = P(A_1)Q\left(\frac{A^*-A_1}{\sigma}\right)+ P(A_2)Q\left(\frac{A_2-A^*}{\sigma}\right)
        \end{equation*}

\item   \begin{enumerate}
        \item En este caso es señalización polar $\pm A$ equiprobable, por lo que el umbral es $0$ y por lo tanto, al usar filtro acoplado tenemos que en recepción:
        \begin{equation*}
                \left(\frac{A}{\sigma}\right)^2_{\max} = 2 \frac{A^2}{\eta} ||p||^2
        \end{equation*}
        Como vimos antes el pulso de Nyquist de exceso $\alpha$ tiene $||p||^2 = T_s(1-\alpha/4) = (3/4)T_s$ por lo que la probabilidad de error es:
        \begin{equation*}
                Q\left(\sqrt{\frac{A}{\sigma}}\right) = Q\left(\sqrt{\frac{3A^2}{2\eta f_s}}\right) \approx Q(6.12) \approx 4.6 \times 10^{-10}. 
        \end{equation*}

        \item El problema del filtro acoplado es que introduce ISI, que en este caso puede ser peor que el ruido. Por lo tanto, podemos usar un filtro pasabajos ideal de ancho $f_s$, el ancho de banda total del pulso. En ese caso, la onda pasará intacta sin ISI y al muestrearla en los lugares correctos obtenemos la densidad de la parte $1$ si las probabilidades son distintas. En este caso, la varianza del ruido es $\sigma^2 = \eta f_s = 10^{-8}$ ya que no es acoplado (hay exceso de ruido).
        
        Calculando para $A_1=-A$, $A_2=A$ y $P(A_1)=\frac{1}{3}=1-P(A_2)$ obtenemos:
        \begin{equation*}
                A^* \approx - 6.93\times 10^{-6}, \quad P_e = 2.7\times 10^{-7}
        \end{equation*}

        El que sean probabilidades diferentes no afecta mucho la probabilidad de error, el problema es el filtro no acoplado. Si fueran equiprobables los símbolos la prob. de error sería $2.86\times 10^{-7}$, bastante mayor que en la parte anterior.
        \end{enumerate}
\end{enumerate}
\end{solucion}

\begin{ejercicio}

Se considera un canal digital de banda base contaminado con ruido blanco de densidad espectral $\eta_0/2 = 10^{-5}\,\unit{\square\volt/\hertz}$. Se utiliza señalización polar con símbolos equiprobables de amplitud $A=3.0\,V$ en la recepción, y velocidad de señalización $f_s=10\,\unit{\kilo\baud}$. Se utilizan pulsos triangulares de la forma:
\begin{equation*}
        p(t)=\begin{cases}
                1-\frac{2|t|}{T_s} &\text{ para } |t|\leq \frac{T_s}{2}\\
                0 & \text{ en otro caso}
            \end{cases}
\end{equation*}
\begin{enumerate}
\item Determinar la mínima probabilidad de error en la detección, posible del sistema. ¿Cómo se obtiene dicho valor?
\item La comunicación anterior se repite en 10 tramos
idénticos, con la amplificación necesaria para restituir el nivel
de señal. Determinar la probabilidad de error a la salida si los
amplificadores: 
\begin{enumerate}
\item no son regenerativos
\item son regenerativos
\end{enumerate}
\end{enumerate}
\end{ejercicio}

\begin{solucion}
        \begin{enumerate}
                \item Si se usa filtro acoplado y señalización polar, la relación a la salida es:
                \begin{equation*}
                        \left(\frac{A}{\sigma}\right)^2_{\max} = 2 \frac{A^2}{\eta}||p||^2
                \end{equation*}
                Para este pulso $||p||^2$ se puede obtener integrando en el tiempo y da $||p||^2 = T_s/3$. Como la señalización es polar, el umbral es $0$ y por lo tanto:
                \begin{equation*}
                        P_e = Q\left(\sqrt{\frac{2A^2}{3\eta f_s}}\right) \approx Q(5.48) \approx 2.2\times 10^{-8}.
                \end{equation*}

                \item \begin{enumerate}
                        \item Si son regenerativos, la probabilidad de error es $\approx m P_e \approx 2.2 \times 10^{-7}$.
                        \item Si no son regeneativos, la $(A/\sigma)^2_{\max}$ se degrada por $m$ por lo que:
                        \begin{equation*}
                                P_e = Q\left(\sqrt{\frac{2A^2}{(10) 3\eta f_s}}\right) \approx Q(1.73) \approx 0.041. 
                        \end{equation*}
                \end{enumerate}
        \end{enumerate}
\end{solucion}


\section{Modulación digital pasabanda}

\begin{ejercicio}
Un sistema de señalización binaria BPSK trasmite información a una velocidad de $2.5 \unit{\mega\bit/\second}$. El ruido en el canal es blanco y gaussiano de densidad espectral de potencia $\eta_0/2=10^{-20}\unit{\square\volt/\hertz}$. La amplitud de la señal a la entrada del detector es $A=1\unit{\micro\volt}$. Asumiendo que se trata de detección coherente con filtro acoplado, determinar la probabilidad de error de bit.
\end{ejercicio}

\begin{solucion}
        La detección coherente BPSK se transforma en banda base a detectar $\pm A$ con umbral $0$ y potencia del ruido $\sigma^2 = \eta f_s$, de donde:
        \begin{equation*}
                P_e = Q\left(\frac{A}{\sigma}\right)=Q\left(\sqrt{\frac{A^2}{\eta f_s}}\right) \approx Q(4.47) \approx 3.9\times 10^{-6}.
        \end{equation*}
\end{solucion}

\begin{ejercicio}
En la detección BPSK coherente, investigar el efecto de un error de fase $\theta$ entre el oscilador local y la portadora. Hallar la expresión para la probabilidad de error de detección.
\end{ejercicio}

\begin{solucion}
Recordemos el diagrama de detección coherente BPSK, pero introduciendo un error de fase en el demodulador:

\begin{blockdiagram}[decoration={brace,amplitude=10}]
        \matrix[row sep=5 mm, column sep=8mm, ampersand replacement=\&]{
            \& \& \& \& \& \node[punto] (p1) {}; \& \node[etiqueta] (muestreo) {};\\
            \node [etiqueta] (p6) {$x_c(t)+n(t)$}; \& \node[operator] (mult) {$\times$}; \& \node [block] (filtro) {$\displaystyle \frac{1}{T_s} \int_{t-T_s}^t$}; \& \node [punto] (p2) {}; \&\node [punto] (p3) {}; \& \&\node [punto] (p4) {}; \&\node [block] (comp) {COMP};\\
            \& \node [etiqueta] (osc) {$2\cos(\omega_ct + \theta)$}; \& \& \& \node [block] (sinc) {SINCR}; \& \node [punto] (p5) {};\\
        };
        \draw [-stealth] (p6.east) --  (mult);
        \draw [-stealth] (mult) -- node [anchor=south] {$x_o(t)$} (filtro);
        \draw [-stealth] (osc) -- (mult);
        %\draw (filtro) -- (p2) -- (p3) -- (muestreo);
        \draw (filtro) -- (p2) -- node [anchor=south] {$y$} (p3) -- (muestreo);
        \draw [-stealth] (p2) |- (sinc);
        \draw (sinc) -- (p5);
        \draw [-stealth,dashed] (p5) -- (p1);
        \draw [-stealth] (p4) -- (comp);
        %\draw [decorate]  (filtro.south east) --node [below=3mm] {\parbox{3cm}{filtro acop. a NRZ, integra en un período de símb.}} (filtro.south west);
    \end{blockdiagram}
    
    con $x_c(t) = \sum_k a_k p(t-kT_s) \cos(\omega_c t)$. Emulando el análisis para detección coherente:
    \begin{align*}
        x_o(t)  &= [x_c(t)+n(t)]2\cos(\omega_c t + \theta) \\
                &= [x_c(t) + n_i(t)\cos(\omega_c t) - n_q(t)\cos(\omega_c t)]2\cos(\omega_c t + \theta)
    \end{align*}
    Aplicando las identidades trigonométricas adecuadas para los términos producto, tendremos:
    \begin{equation*}
        x_o(t) = \sum_k a_k p(t-kT_s) \cos(\theta) + n_i(t)\cos(\theta) - n_q(t)\sin(\theta) + \ldots
    \end{equation*}
    donde los puntos representan términos de alta frecuencia que se irán luego de filtrar y ya los despreciamos.
    
    Como los ruidos $n_i(t)$ y $n_q(t)$ son independientes, la combinación lineal dará un ruido $n_b(t)$ en banda base de densidad espectral $\eta\cos^2(\theta) + \eta \sin^2(\theta) = \eta$, por lo que la potencia del ruido no cambia por el error de fase.

    En cambio, la amplitud del símbolo $a_k$ se ve atenuada por un factor $\cos(\theta)$ que depende del error. Si utilizamos el filtro acoplado, lograremos:
    \begin{equation*}
        \left( \frac{A'}{\sigma}\right)^2_{\max} = \frac{A^2 \cos^2(\theta)}{\eta}||p||^2 = \frac{A^2 \cos^2(\theta)}{\eta f_s},
    \end{equation*}
    de donde
    \begin{equation*}
        Q\left(\sqrt{\frac{A'}{\sigma}}\right) = Q\left(\sqrt{\frac{A^2 \cos^2(\theta)}{\eta f_s}}\right).
    \end{equation*}

    En particular, si $\theta = \pm \pi/2$ estoy completamente fuera de fase y solo queda el ruido, pierdo completamente información cruzada del símbolo (prob. de error $\to 1/2$). Si en cambio $\theta>\pi/2$ tenemos inversión de símbolos, porque cambia el signo de $\cos(\theta)$. En el caso extremo de $\theta=\pi$ (medio ciclo), decodificamos el mensaje al revés.

\end{solucion}

\begin{ejercicio}
Considere el sistema QPSK cuya señal transmitida es
\begin{equation*}
x_c(t)= \sum_k A p_R(t-kT_s)\cos(2\pi f_c t + \theta_k ),\quad\quad \theta_k\in\left\{0,\frac{\pi}{2},\pi,\frac{3\pi}{2}\right\}.
\end{equation*}
Aquí $p_R(t)$ es el pulso rectangular NRZ.
\begin{enumerate}
\item Se dispone de dos generadores locales, $\cos(\omega_c t+\pi/4)$ y $\cos(\omega_c t+3\pi/4)$. Diseñar en diagrama de bloques un detector sincrónico
apropiado usando filtros acoplados.
\item Para el detector diseñado, calcular la probabilidad de error de símbolo y expresarla en función del parámetro
$\gamma:=\frac{\overline{E}_p}{\eta}$.
\end{enumerate}
\end{ejercicio}

\begin{solucion}
        \begin{enumerate}
        \item  Si multiplicamos $x_c$ por el generador local $2\cos(\omega_c t + \pi/4)$ y despreciamos los términos de alta frecuencia, nos queda:
        \begin{align*}
                x_1(t) &= \sum_k Ap_R(t-kT_s)\cos(\omega_c t + \theta_k)2\cos(\omega_c t + \pi/4) \\
                &= \sum_k Ap_R(t-kT_s)\cos(\theta_k-\pi/4) + \ldots
        \end{align*}
        Con el mismo razonamiento, multiplicando por el otro oscilador:
        \begin{align*}
                x_2(t) &= \sum_k Ap_R(t-kT_s)\cos(\omega_c t + \theta_k)2\cos(\omega_c t + 3\pi/4) \\
                &= \sum_k Ap_R(t-kT_s)\cos(\theta_k-3\pi/4) + \ldots
        \end{align*}

        Notemos que si $\theta_k=0$ o $\pi/2$, $x_1$ produce la misma señal de amplitud $(A/\sqrt{2})p_R$, pero $x_2$ cambia de signo. Razonando de este modo, hacemos un filtro acoplado a cada una de las señales y regeneramos de acuerdo a:
        \begin{equation*}
                \begin{cases}
                        x_1 = A/\sqrt{2}, x_2 = -A/\sqrt{2} & \mapsto \theta_k = 0,\\
                        x_1 = A/\sqrt{2}, x_2 = A/\sqrt{2} & \mapsto \theta_k = \pi/2,\\
                        x_1 = -A/\sqrt{2}, x_2 = A/\sqrt{2} & \mapsto \theta_k = \pi,\\
                        x_1 = -A/\sqrt{2}, x_2 = -A/\sqrt{2} & \mapsto \theta_k = 3\pi/4.
                \end{cases}
        \end{equation*}

        \item Para el sistema anterior, la probabilidad de error de símbolo es igual a un sistema $4-QAM$ con amplitud $A/\sqrt{2}$. En cada señal usamos $0$ como umbral por lo que la probabilidad de error en $x_1$ o en $x_2$ es $Q(A/(\sqrt{2}\sigma))$ y la probabilidad total es $P_e\approx 2Q(A/(\sqrt{2}\sigma))$.
        
        Observando que $\bar{E}_p = A^2 T_s/2$ y $\sigma^2=\eta f_s$, la expresión queda $P_e \approx 2Q\left(\sqrt{\frac{\bar{E}_p}{\eta}}\right) = 2Q(\sqrt{\gamma})$.
\end{enumerate}
\end{solucion}

\begin{ejercicio}
Se considera un sistema de transmisión digital pasabanda que utiliza los siguientes $M=8$ pulsos de radio frecuencia 
\begin{align*} & A p_R(t)[\pm \cos(2\pi f_c t) \pm \sen(2\pi f_c t)],\\ & A p_R(t)[\pm 3\cos(2\pi f_c t)
\pm \sen(2\pi f_c t)].
\end{align*}
Aquí $p_R(t)$ es el pulso rectangular NRZ. 

\begin{enumerate}
\item Dibujar un diagrama de la constelación de símbolos
utilizada. Mostrar una codificación de Gray para representar 3
bits mediante esta constelación.
\item Suponiendo símbolos equiprobables, hallar la densidad
espectral de potencia de la señal digital correspondiente.
\item Diseñar en diagrama de bloques un detector apropiado para este sistema.
\item Si el canal presenta ruido blanco gaussiano de densidad espectral de
potencia $\eta/2$, hallar una fórmula para la probabilidad $P_e$
de error de símbolo para el sistema, en función de
$\gamma:=\frac{\bar{E}_p}{\eta}$, siendo $\bar{E}_p$ la energía
media de los pulsos.
\item Evaluar numéricamente la curva $P_e(\gamma)$ de (d), y comparar con
la correspondiente a modulación 8PSK. ?`Cuál es mejor?
\item Comparar también con la curva para una transmisión de banda
base polar, con $M=8$ símbolos. ?`Es justa esta comparación en
términos de eficiencia espectral?
\end{enumerate}
\end{ejercicio}

\begin{solucion}
        \begin{enumerate}
                \item Tomando las componentes que multiplican a $\sin$ y $\cos$ como componentes en fase y cuadratura en banda base, los símbolos complejos son:
                \begin{center}
                        \begin{tikzpicture}
                                \begin{axis}[x=1cm, y=1cm, axis lines=middle,  
                                        xlabel=$a_k$,
                                        ylabel=$b_k$,
                                        x label style={anchor=north west},
                                        y label style={anchor=south east},
                                        xmin=-3.5,xmax=3.5,ymin=-1.5,ymax=1.5,
                                        xtick={-3,-1,0,1,3}, ytick={-1,0,1},
                                        xticklabels={$-3A$, $-A$, $0$, $A$, $3A$},
                                        yticklabels={$-A$, $0$, $A$},
                                        clip=false,
                                    ]
                                \addplot+[only marks, teal, mark options={fill=teal}] coordinates {
                                        (-3,-1) (-3,1) (-1,-1) (-1,1) 
                                        (1,-1) (1,1) (3,-1) (3,1)
                                };
                                \node[below] at (axis cs:-3,-1) {\tiny (000)};
                                \node[below] at (axis cs:-3,1) {\tiny (100)};
                                \node[below] at (axis cs:-1,-1) {\tiny (001)};
                                \node[below] at (axis cs:-1,1) {\tiny (101)};
                                \node[below] at (axis cs:1,-1) {\tiny (011)};
                                \node[below] at (axis cs:1,1) {\tiny (111)};
                                \node[below] at (axis cs:3,-1) {\tiny (010)};
                                \node[below] at (axis cs:3,1) {\tiny (110)};

                                \end{axis}
                        \end{tikzpicture}
                \end{center}

                Ya agregamos la codificación de Gray adecuada (básicamente dos codificaciones de Gray de 4 bits para $a_k$, y el bit mas significativo determina el $b_k$).

                \item Observando que la constelación está centrada en $0$, $\mu=0$ y calculando los módulos de la constelación:
                      \begin{equation*}
                      E[c_k^2] = E[a_k^2+b_k^2] = 4 (1/8) (A^2 + A^2) + 4 (1/8) (A^2+9A^2) = 6A^2
                      \end{equation*}

                La densidad espectral de potencia en banda base será:
                \begin{equation*}
                        S_{x_{ce}} (f) = 6A^2 f_s |\hat{P}(f)|^2 = 6A^2 T_s \sinc^2(\pi T_s f).
                \end{equation*}
                donde se usó que el pulso es NRZ. Luego se modula obteniendo:
                \begin{equation*}
                        S_{x_c} = \frac{1}{4}\left[S_{x_{ce}}(f-f_c) + S_{x_{ce}}(f+f_c)\right].
                \end{equation*}

                Bosquejo (para $f>0$):
                \begin{center}
                        \begin{tikzpicture}
                        \pgfmathdeclarefunction{sinc}{1}{%
                        \pgfmathparse{(#1==0 ? 1: sin(#1 r))/(#1==0 ? 1: #1)}%
                        }
                                \begin{axis}[x=1cm, y=1cm, axis lines=middle,  
                                        xlabel=$f$,
                                        ylabel=$S_{x_c}(f)$,
                                        x label style={anchor=north west},
                                        y label style={anchor=south east},
                                        xmin=-.1,xmax=8,ymin=-.1,ymax=1.5,
                                        xtick={3.5,5,6.5}, ytick=\empty,
                                        xticklabels={$f_c-f_s$,$f_c$,$f_c+f_s$},
                                        clip=false,
                                    ]
                                \addplot+[solid, mark=none, very thick, teal, domain = (2:8), samples=100]  {
                                        sinc(pi*(x-5)/1.5)^2
                                } node[midway, above right] {$\frac{3}{2}A^2 T_s$};

                                \draw[dashed] (axis cs:5,0) -- (axis cs:5,1.5);
                                \end{axis}

                        \end{tikzpicture}
                \end{center}

                \item Podemos aplicar detección coherente multiplicando por $2\cos(\omega_c t)$ y $-2\sin(\omega_c t)$ y luego aplicando filtro acoplado. Esto es equivalente a:
                 \begin{itemize}
                        \item Una detección $M-$aria para $a_k$ en banda base, con $M=4$.
                        \item Una detección binaria para $b_k$.
                 \end{itemize}

                 \item Usamos la observación anterior para la probabilidad de error:
                     \begin{align*}
                        P_e(a_k) &= 2\frac{M-1}{M}Q(A/\sigma) = \frac{3}{2}Q(A/\sigma),\\
                        P_e(b_k) &= Q(A/\sigma),\\
                        P_e &\approx P_e(a_k) + P_e(b_k) = \frac{5}{2}Q(A/\sigma).
                     \end{align*}

                     Como el filtro es acoplado, se tiene que:
                     \begin{equation*}
                        \left(\frac{A}{\sigma}\right)^2_{\max} = 2 \frac{A^2}{\eta}||p||^2 = 2\frac{A^2}{\eta f_s}.
                     \end{equation*}

                     Si lo llevamos a energía media del pulso de radiofrecuencia, se tiene que:
                     \begin{equation*}
                        \bar{E}_p = E[c_k^2]||\tilde{p}||^2 = 6 A^2 T_s/2 = 3A^2T_s \Rightarrow \frac{A^2}{f_s} = \frac{\bar{E}_p}{3}
                     \end{equation*}
                     De donde:
                     \begin{equation*}
                        \left(\frac{A}{\sigma}\right)^2_{\max} = \frac{2\bar{E}_p}{3\eta f_s} \Rightarrow P_e = Q\left(\sqrt{\frac{2\bar{E}_p}{3\eta f_s} }\right) = Q\left(\sqrt{2\gamma/3}\right).
                     \end{equation*}
                

                \item Se tiene la siguiente comparación, usando que la probabilidad de error de $8-$PSK es:
                \begin{equation*}
                        P_e \approx 2 Q\left(\sqrt{2\gamma\sin^2(\pi/M)}\right)
                \end{equation*}
                \begin{center}
                        \begin{tikzpicture}

\begin{axis}[   width = 0.8\textwidth, height = 0.5\textwidth,
                xlabel=$\gamma$ (dB),
                ylabel=$P_e$,
                x label style={anchor=north west},
                y label style={anchor=south east},
                xmin=3,xmax=28,ymin=-20,ymax=0, grid,
                ytick = {-20,-15,-10,-5,0},
                yticklabels = {$10^{-20}$,$10^{-15}$,$10^{-10}$,$10^{-5}$,$10^0$},
            ]
    \addplot[color=teal, very thick, mark=none, solid]
        table[row sep={\\}]
        {
            \\
            3.0  -0.9052248154440873  \\
            3.2727272727272725  -0.9318155274977596  \\
            3.5454545454545454  -0.9597607678454317  \\
            3.8181818181818183  -0.9891390862577611  \\
            4.090909090909091  -1.0200340060905306  \\
            4.363636363636363  -1.052534352874508  \\
            4.636363636363637  -1.0867346046994981  \\
            4.909090909090909  -1.1227352658019767  \\
            5.181818181818182  -1.1606432648536253  \\
            5.454545454545454  -1.2005723795415277  \\
            5.7272727272727275  -1.242643689130117  \\
            6.0  -1.2869860568005904  \\
            6.272727272727273  -1.3337366436759102  \\
            6.545454545454545  -1.3830414565591027  \\
            6.8181818181818175  -1.43505593153997  \\
            7.090909090909091  -1.4899455557609873  \\
            7.363636363636363  -1.5478865297777458  \\
            7.636363636363637  -1.609066473103435  \\
            7.909090909090909  -1.6736851756911681  \\
            8.181818181818182  -1.7419553982832692  \\
            8.454545454545455  -1.8141037247435903  \\
            8.727272727272727  -1.8903714696884664  \\
            9.0  -1.9710156449448093  \\
            9.272727272727273  -2.0563099885910403  \\
            9.545454545454547  -2.146546060579076  \\
            9.818181818181818  -2.2420344091943316  \\
            10.09090909090909  -2.343105812886911  \\
            10.363636363636363  -2.4501126023018034  \\
            10.636363636363637  -2.5634300676503563  \\
            10.909090909090908  -2.683457956900691  \\
            11.181818181818182  -2.810622070622481  \\
            11.454545454545455  -2.945375959703007  \\
            11.727272727272727  -3.0882027325582007  \\
            12.0  -3.2396169788959535  \\
            12.272727272727273  -3.4001668175510713  \\
            12.545454545454547  -3.5704360764037317  \\
            12.818181818181817  -3.751046612917804  \\
            13.09090909090909  -3.942660784394253  \\
            13.363636363636363  -4.145984077629879  \\
            13.636363636363635  -4.361767908305417  \\
            13.909090909090908  -4.5908126011016686  \\
            14.181818181818182  -4.833970562260641  \\
            14.454545454545455  -5.092149657073337  \\
            14.727272727272727  -5.366316805589741  \\
            15.0  -5.6575018107128265  \\
            15.272727272727273  -5.966801433760568  \\
            15.545454545454545  -6.295383733561263  \\
            15.818181818181818  -6.644492686192081  \\
            16.09090909090909  -7.0154531035824785  \\
            16.363636363636363  -7.409675870387399  \\
            16.636363636363637  -7.828663519794765  \\
            16.90909090909091  -8.274016170272324  \\
            17.181818181818183  -8.747437846685962  \\
            17.454545454545453  -9.250743210740731  \\
            17.727272727272727  -9.785864727312914  \\
            18.0  -10.354860294962913  \\
            18.272727272727273  -10.959921370751578  \\
            18.545454545454547  -11.603381621433696  \\
            18.81818181818182  -12.287726135179913  \\
            19.090909090909093  -13.015601230189851  \\
            19.363636363636363  -13.789824898914521  \\
            19.636363636363637  -14.613397929113257  \\
            19.90909090909091  -15.489515745640395  \\
            20.18181818181818  -16.421581019699666  \\
            20.454545454545453  -17.413217095331266  \\
            20.727272727272727  -18.468282286119795  \\
            21.0  -19.590885098543254  \\
        }
        ;
    \addlegendentry {Mod. propuesta}
    \addplot[color=red, very thick, mark=none, solid]
        table[row sep={\\}]
        {
            \\
            3.0  -0.3520371079914345  \\
            3.2727272727272725  -0.36632626225770565  \\
            3.5454545454545454  -0.38128083493782927  \\
            3.8181818181818183  -0.3969366070276897  \\
            4.090909090909091  -0.4133315219723772  \\
            4.363636363636363  -0.43050582638930673  \\
            4.636363636363637  -0.448502220247813  \\
            4.909090909090909  -0.4673660171374084  \\
            5.181818181818182  -0.4871453152975161  \\
            5.454545454545454  -0.5078911801245576  \\
            5.7272727272727275  -0.5296578389179419  \\
            6.0  -0.5525028886749117  \\
            6.272727272727273  -0.5764875177955255  \\
            6.545454545454545  -0.6016767426134648  \\
            6.8181818181818175  -0.6281396597260421  \\
            7.090909090909091  -0.6559497151579643  \\
            7.363636363636363  -0.6851849914582836  \\
            7.636363636363637  -0.7159285138987946  \\
            7.909090909090909  -0.7482685770151685  \\
            8.181818181818182  -0.7822990928096305  \\
            8.454545454545455  -0.8181199620162912  \\
            8.727272727272727  -0.8558374699176586  \\
            9.0  -0.8955647082937459  \\
            9.272727272727273  -0.937422025183929  \\
            9.545454545454547  -0.981537504246704  \\
            9.818181818181818  -1.0280474756142106  \\
            10.09090909090909  -1.0770970602572685  \\
            10.363636363636363  -1.1288407500032887  \\
            10.636363636363637  -1.1834430254842414  \\
            10.909090909090908  -1.2410790144355748  \\
            11.181818181818182  -1.3019351929201306  \\
            11.454545454545455  -1.3662101322144133  \\
            11.727272727272727  -1.434115294268753  \\
            12.0  -1.5058758788387232  \\
            12.272727272727273  -1.5817317255834118  \\
            12.545454545454547  -1.6619382746377227  \\
            12.818181818181817  -1.7467675893916432  \\
            13.09090909090909  -1.8365094454504003  \\
            13.363636363636363  -1.9314724900065263  \\
            13.636363636363635  -2.031985476129347  \\
            13.909090909090908  -2.1383985767701947  \\
            14.181818181818182  -2.2510847835940946  \\
            14.454545454545455  -2.370441396082017  \\
            14.727272727272727  -2.4968916067032616  \\
            15.0  -2.630886188336701  \\
            15.272727272727273  -2.7729052905238336  \\
            15.545454545454545  -2.9234603515675164  \\
            15.818181818181818  -3.083096133949505  \\
            16.09090909090909  -3.2523928910293174  \\
            16.363636363636363  -3.4319686735082997  \\
            16.636363636363637  -3.622481784698155  \\
            16.90909090909091  -3.824633394224657  \\
            17.181818181818183  -4.039170320427079  \\
            17.454545454545453  -4.266887992384542  \\
            17.727272727272727  -4.508633603214359  \\
            18.0  -4.765309467047453  \\
            18.272727272727273  -5.037876592894983  \\
            18.545454545454547  -5.327358489481303  \\
            18.81818181818182  -5.634845216034938  \\
            19.090909090909093  -5.961497695004684  \\
            19.363636363636363  -6.30855230370622  \\
            19.636363636363637  -6.677325763009611  \\
            19.90909090909091  -7.069220342354187  \\
            20.18181818181818  -7.485729401629287  \\
            20.454545454545453  -7.928443291791732  \\
            20.727272727272727  -8.399055637509344  \\
            21.0  -8.899370026629711  \\
            21.272727272727273  -9.431307132880702  \\
            21.545454545454547  -9.99691229992016  \\
            21.818181818181817  -10.598363616674053  \\
            22.09090909090909  -11.237980515841485  \\
            22.363636363636363  -11.91823292950975  \\
            22.636363636363637  -12.641751038021148  \\
            22.90909090909091  -13.411335650573637  \\
            23.181818181818183  -14.229969258529758  \\
            23.454545454545453  -15.100827805061728  \\
            23.727272727272727  -16.027293217586248  \\
            24.0  -17.01296675245083  \\
            24.272727272727273  -18.06168320453733  \\
            24.545454545454547  -19.17752603785947  \\
        }
        ;
    \addlegendentry {8-PSK}
\end{axis}
\end{tikzpicture}

                \end{center}

                \item Se tiene la siguiente comparación, usando la expresión para un sistema $M-$ario en banda base:
                \begin{equation*}
                        P_e = 2 \frac{M-1}{M} Q\left(\sqrt{\frac{3\gamma}{M^2-1}}\right)
                \end{equation*}
                \begin{center}
                        \begin{tikzpicture}

\begin{axis}[   width = 0.8\textwidth, height = 0.5\textwidth,
                xlabel=$\gamma$ (dB),
                ylabel=$P_e$,
                x label style={anchor=north west},
                y label style={anchor=south east},
                xmin=3,xmax=28,ymin=-20,ymax=0, grid,
                clip=true,
                ytick = {-20,-15,-10,-5,0},
                yticklabels = {$10^{-20}$,$10^{-15}$,$10^{-10}$,$10^{-5}$,$10^0$},
            ]
    \addplot[color=teal, very thick, mark=none, solid]
        table[row sep={\\}]
        {
            \\
            3.0  -0.9052248154440873  \\
            3.2727272727272725  -0.9318155274977596  \\
            3.5454545454545454  -0.9597607678454317  \\
            3.8181818181818183  -0.9891390862577611  \\
            4.090909090909091  -1.0200340060905306  \\
            4.363636363636363  -1.052534352874508  \\
            4.636363636363637  -1.0867346046994981  \\
            4.909090909090909  -1.1227352658019767  \\
            5.181818181818182  -1.1606432648536253  \\
            5.454545454545454  -1.2005723795415277  \\
            5.7272727272727275  -1.242643689130117  \\
            6.0  -1.2869860568005904  \\
            6.272727272727273  -1.3337366436759102  \\
            6.545454545454545  -1.3830414565591027  \\
            6.8181818181818175  -1.43505593153997  \\
            7.090909090909091  -1.4899455557609873  \\
            7.363636363636363  -1.5478865297777458  \\
            7.636363636363637  -1.609066473103435  \\
            7.909090909090909  -1.6736851756911681  \\
            8.181818181818182  -1.7419553982832692  \\
            8.454545454545455  -1.8141037247435903  \\
            8.727272727272727  -1.8903714696884664  \\
            9.0  -1.9710156449448093  \\
            9.272727272727273  -2.0563099885910403  \\
            9.545454545454547  -2.146546060579076  \\
            9.818181818181818  -2.2420344091943316  \\
            10.09090909090909  -2.343105812886911  \\
            10.363636363636363  -2.4501126023018034  \\
            10.636363636363637  -2.5634300676503563  \\
            10.909090909090908  -2.683457956900691  \\
            11.181818181818182  -2.810622070622481  \\
            11.454545454545455  -2.945375959703007  \\
            11.727272727272727  -3.0882027325582007  \\
            12.0  -3.2396169788959535  \\
            12.272727272727273  -3.4001668175510713  \\
            12.545454545454547  -3.5704360764037317  \\
            12.818181818181817  -3.751046612917804  \\
            13.09090909090909  -3.942660784394253  \\
            13.363636363636363  -4.145984077629879  \\
            13.636363636363635  -4.361767908305417  \\
            13.909090909090908  -4.5908126011016686  \\
            14.181818181818182  -4.833970562260641  \\
            14.454545454545455  -5.092149657073337  \\
            14.727272727272727  -5.366316805589741  \\
            15.0  -5.6575018107128265  \\
            15.272727272727273  -5.966801433760568  \\
            15.545454545454545  -6.295383733561263  \\
            15.818181818181818  -6.644492686192081  \\
            16.09090909090909  -7.0154531035824785  \\
            16.363636363636363  -7.409675870387399  \\
            16.636363636363637  -7.828663519794765  \\
            16.90909090909091  -8.274016170272324  \\
            17.181818181818183  -8.747437846685962  \\
            17.454545454545453  -9.250743210740731  \\
            17.727272727272727  -9.785864727312914  \\
            18.0  -10.354860294962913  \\
            18.272727272727273  -10.959921370751578  \\
            18.545454545454547  -11.603381621433696  \\
            18.81818181818182  -12.287726135179913  \\
            19.090909090909093  -13.015601230189851  \\
            19.363636363636363  -13.789824898914521  \\
            19.636363636363637  -14.613397929113257  \\
            19.90909090909091  -15.489515745640395  \\
            20.18181818181818  -16.421581019699666  \\
            20.454545454545453  -17.413217095331266  \\
            20.727272727272727  -18.468282286119795  \\
            21.0  -19.590885098543254  \\
        }
        ;
    \addlegendentry {Mod. propuesta}
    \addplot[color=red, very thick, mark=none, solid]
        table[row sep={\\}]
        {
            \\
            3.0  -0.2365466627224427  \\
            3.2727272727272725  -0.24318650134922756  \\
            3.5454545454545454  -0.2501017951330155  \\
            3.8181818181818183  -0.25730568116263286  \\
            4.090909090909091  -0.2648120174938776  \\
            4.363636363636363  -0.27263542739851565  \\
            4.636363636363637  -0.28079134653285154  \\
            4.909090909090909  -0.2892960732254347  \\
            5.181818181818182  -0.29816682209710443  \\
            5.454545454545454  -0.3074217812411255  \\
            5.7272727272727275  -0.3170801732066669  \\
            6.0  -0.32716232004537943  \\
            6.272727272727273  -0.3376897126984039  \\
            6.545454545454545  -0.3486850850198517  \\
            6.8181818181818175  -0.36017249275269847  \\
            7.090909090909091  -0.37217739779420145  \\
            7.363636363636363  -0.38472675811046364  \\
            7.636363636363637  -0.397849123683685  \\
            7.909090909090909  -0.41157473890106666  \\
            8.181818181818182  -0.42593565182132986  \\
            8.454545454545455  -0.440965830783488  \\
            8.727272727272727  -0.4567012888529386  \\
            9.0  -0.47318021663225235  \\
            9.272727272727273  -0.4904431239982991  \\
            9.545454545454547  -0.5085329913637128  \\
            9.818181818181818  -0.52749543109925  \\
            10.09090909090909  -0.5473788597944815  \\
            10.363636363636363  -0.5682346820776151  \\
            10.636363636363637  -0.590117486761193  \\
            10.909090909090908  -0.6130852561291461  \\
            11.181818181818182  -0.6371995892323328  \\
            11.454545454545455  -0.6625259401144512  \\
            11.727272727272727  -0.6891338719482839  \\
            12.0  -0.7170973281238022  \\
            12.272727272727273  -0.7464949213949609  \\
            12.545454545454547  -0.7774102422612944  \\
            12.818181818181817  -0.8099321878339314  \\
            13.09090909090909  -0.8441553125136897  \\
            13.363636363636363  -0.8801802018917357  \\
            13.636363636363635  -0.9181138713713312  \\
            13.909090909090908  -0.9580701911026809  \\
            14.181818181818182  -1.0001703389223013  \\
            14.454545454545455  -1.0445432830940702  \\
            14.727272727272727  -1.0913262967615678  \\
            15.0  -1.1406655061410613  \\
            15.272727272727273  -1.1927164746119443  \\
            15.545454545454545  -1.2476448249972436  \\
            15.818181818181818  -1.3056269024715146  \\
            16.09090909090909  -1.3668504806876745  \\
            16.363636363636363  -1.4315155138788318  \\
            16.636363636363637  -1.4998349378665572  \\
            16.90909090909091  -1.5720355230942322  \\
            17.181818181818183  -1.6483587830037365  \\
            17.454545454545453  -1.7290619412868793  \\
            17.727272727272727  -1.8144189617703417  \\
            18.0  -1.9047216449355966  \\
            18.272727272727273  -2.0002807953343007  \\
            18.545454545454547  -2.101427464436019  \\
            18.81818181818182  -2.2085142737401258  \\
            19.090909090909093  -2.3219168232983556  \\
            19.363636363636363  -2.4420351911302256  \\
            19.636363636363637  -2.5692955293715753  \\
            19.90909090909091  -2.7041517633782846  \\
            20.18181818181818  -2.847087400414374  \\
            20.454545454545453  -2.9986174549876177  \\
            20.727272727272727  -3.159290498358291  \\
            21.0  -3.3296908402395493  \\
            21.272727272727273  -3.510440851232876  \\
            21.545454545454547  -3.7022034351013633  \\
            21.818181818181817  -3.905684660579121  \\
            22.09090909090909  -4.1216365630493685  \\
            22.363636363636363  -4.350860127098968  \\
            22.636363636363637  -4.594208461676121  \\
            22.90909090909091  -4.852590180343145  \\
            23.181818181818183  -5.126972999930836  \\
            23.454545454545453  -5.418387571767978  \\
            23.727272727272727  -5.727931560582415  \\
            24.0  -6.056773987152154  \\
            24.272727272727273  -6.406159851830627  \\
            24.545454545454547  -6.777415057182635  \\
            24.81818181818182  -7.171951649151939  \\
            25.090909090909093  -7.591273397441952  \\
            25.36363636363636  -8.036981737132644  \\
            25.636363636363633  -8.510782094985178  \\
            25.909090909090907  -9.01449062540576  \\
            26.18181818181818  -9.550041382659085  \\
            26.454545454545453  -10.119493957644295  \\
            26.727272727272727  -10.725041609380726  \\
            27.0  -11.36901992330363  \\
            27.27272727272727  -12.05391603054899  \\
            27.545454545454543  -12.782378424620266  \\
            27.818181818181817  -13.557227414186686  \\
            28.09090909090909  -14.381466253272325  \\
            28.363636363636363  -15.258292992767158  \\
            28.636363636363637  -16.1911131000364  \\
            28.90909090909091  -17.183552896433934  \\
            29.18181818181818  -18.239473865751503  \\
            29.454545454545453  -19.362987890070293  \\
        }
        ;
    \addlegendentry {$M$-ario Banda base}
\end{axis}
\end{tikzpicture}

                \end{center}
        \end{enumerate}
\end{solucion}


\begin{ejercicio}
En los sistemas coherentes donde se debe enviar una pequeña señal piloto para la sincronización, empeora la eficiencia. Considere un sistema digital BPSK, cuya onda trasmitida es
\begin{equation*}
        x_c(t) = \sum_k a_k\sqrt{1-\mu^2}p_R(t-kT_s)\cos(\omega_c t)+\mu
        A\sin(\omega_c t),
\end{equation*}
donde $a_k=\{+A,-A\}$, $0<\mu<1$, y $p_R(t)$ es el pulso rectangular NRZ de largo $T_s$. Calcule la probabilidad de error de símbolo en función de la densidad espectral de potencia del ruido $\frac{\eta}{2}$ y la potencia media trasmitida. Compare el resultado con el que se obtendría si no se enviara la portadora piloto $s_c(t)=\mu A\sin(\omega_c t)$.
\end{ejercicio}

\begin{solucion}
La onda enviada tiene una componente BPSK y una componente de portadora. La componente BPSK tiene amplitudes $a'_k = \pm A\sqrt{1-\mu^2}$ por lo que la potencia es $S_{\mathrm{bpsk}} = E[(a'_k)^2] f_s ||\tilde{p}||^2$ siendo $||\tilde{p}||$ el pulso NRZ en radiofrecuencia, de energía $T_s/2$. Por lo tanto: 
\begin{equation*}
        S_{\mathrm{bpsk}}=\frac{1}{2}E[(a'_k)^2] = \frac{A^2 (1-\mu^2)}{2}.        
\end{equation*}

La parte sinusoidal de portadora es de potencia $\frac{A^2\mu^2}{2}$ por lo que la potencia total transmitida es $S=A^2$. Es decir, usamos la misma potencia total que un BPSK común, pero usamos un índice de modulación $\mu$ para la portadora, reservando una fracción de potencia para la misma.

Si aplicamos detección coherente a la parte BPSK, tenemos que usar $A' = A\sqrt{1-\mu^2}$ por lo que sustituyendo en la expresión de BPSK queda:
\begin{equation*}
        P_e = Q\left(A'/\sigma\right) = Q\left(\sqrt{\frac{A^2(1-\mu^2)}{\eta f_s}}\right) = Q\left(\sqrt{2\gamma(1-\mu^2)}\right),
\end{equation*}
donde usamos que $\gamma = \bar{E}_p/\eta$ con $\bar{E}_p = A^2/2$, la energía del pulso equivalente en BPSK sin portadora. A continuación una comparación del efecto de $\mu$ para $\gamma = 10$ dB.
\begin{center}
        
\begin{tikzpicture}

\begin{axis}[   width = 0.8\textwidth, height = 0.5\textwidth,
                xlabel=$\mu$,
                ylabel=$P_e$,
                x label style={anchor=north west},
                y label style={anchor=south east},
                xmin=0,xmax=1,ymin=-7,ymax=0, grid,
                ytick = {-6,-5,-4,-3,-2,-1},
                yticklabels = {$10^{-6}$,$10^{-5}$,$10^{-4}$,$10^{-3}$,$10^{-2}$, $10^{-1}$},
                clip=true,
            ]

    \addplot[color=teal, very thick, mark=none, solid]
        table[row sep={\\}]
        {   \\
            0.0  -5.412052513773938  \\
            0.01  -5.411598271072702  \\
            0.02  -5.410235531873896  \\
            0.03  -5.4079642628714835  \\
            0.04  -5.404784408486221  \\
            0.05  -5.400695890761566  \\
            0.06  -5.3956986092173995  \\
            0.07  -5.3897924406609485  \\
            0.08  -5.382977238954171  \\
            0.09  -5.375252834736659  \\
            0.1  -5.3666190351029535  \\
            0.11  -5.35707562323295  \\
            0.12  -5.346622357973901  \\
            0.13  -5.335258973372275  \\
            0.14  -5.322985178153527  \\
            0.15  -5.309800655147599  \\
            0.16  -5.295705060657705  \\
            0.17  -5.280698023769695  \\
            0.18  -5.26477914559895  \\
            0.19  -5.247947998471547  \\
            0.2  -5.230204125035956  \\
            0.21  -5.211547037301272  \\
            0.22  -5.191976215597525  \\
            0.23  -5.171491107453153  \\
            0.24  -5.150091126384317  \\
            0.25  -5.1277756505901095  \\
            0.26  -5.104544021547194  \\
            0.27  -5.080395542496766  \\
            0.28  -5.05532947681598  \\
            0.29  -5.029345046265278  \\
            0.3  -5.002441429102125  \\
            0.31  -4.974617758050746  \\
            0.32  -4.945873118116396  \\
            0.33  -4.9162065442315015  \\
            0.34  -4.885617018719695  \\
            0.35  -4.854103468562346  \\
            0.36  -4.821664762450483  \\
            0.37  -4.788299707603228  \\
            0.38  -4.754007046331764  \\
            0.39  -4.71878545232557  \\
            0.4  -4.682633526635028  \\
            0.41  -4.645549793321553  \\
            0.42  -4.607532694743137  \\
            0.43  -4.568580586439291  \\
            0.44  -4.528691731575229  \\
            0.45  -4.487864294900158  \\
            0.46  -4.4460963361690045  \\
            0.47  -4.403385802970566  \\
            0.48  -4.35973052289775  \\
            0.49  -4.315128194987208  \\
            0.5  -4.2695763803460185  \\
            0.51  -4.223072491871933  \\
            0.52  -4.175613782960816  \\
            0.53  -4.127197335079887  \\
            0.54  -4.077820044068  \\
            0.55  -4.027478605003746  \\
            0.56  -3.9761694954583584  \\
            0.57  -3.923888956922274  \\
            0.58  -3.8706329741611247  \\
            0.59  -3.8163972522176945  \\
            0.6  -3.761177190729801  \\
            0.61  -3.7049678551784417  \\
            0.62  -3.64776394461392  \\
            0.63  -3.589559755327531  \\
            0.64  -3.5303491398394873  \\
            0.65  -3.4701254604560785  \\
            0.66  -3.408881536505456  \\
            0.67  -3.346609584185184  \\
            0.68  -3.28330114773723  \\
            0.69  -3.218947020396183  \\
            0.7  -3.1535371532193737  \\
            0.71  -3.087060549483785  \\
            0.72  -3.0195051417980614  \\
            0.73  -2.9508576483935496  \\
            0.74  -2.8811034041782837  \\
            0.75  -2.8102261609967356  \\
            0.76  -2.738207850044808  \\
            0.77  -2.665028297416139  \\
            0.78  -2.5906648811201327  \\
            0.79  -2.515092114351748  \\
            0.8  -2.438281134923716  \\
            0.81  -2.3601990740220185  \\
            0.82  -2.2808082679507606  \\
            0.83  -2.2000652629600883  \\
            0.84  -2.1179195435023908  \\
            0.85  -2.0343118849579516  \\
            0.86  -1.9491721874277297  \\
            0.87  -1.8624165781156035  \\
            0.88  -1.7739434594552985  \\
            0.89  -1.6836279981584374  \\
            0.9  -1.5913142392544668  \\
            0.91  -1.4968034744968095  \\
            0.92  -1.3998364552567186  \\
            0.93  -1.3000649738554546  \\
            0.94  -1.1970039202082758  \\
            0.95  -1.0899445767177687  \\
            0.96  -0.9777826406644201  \\
            0.97  -0.8586298632249533  \\
            0.98  -0.728743978080803  \\
            0.99  -0.5782941313571539  \\
            1.0  -0.3010299956639812  \\
        };
\end{axis}
\end{tikzpicture}

\end{center}

\end{solucion}

\begin{ejercicio}
El sistema de la figura es un detector diferencial no coherente para
señales BPSK.

\begin{center}
\begin{tikzpicture}
        [node distance=2cm,>=stealth, auto,
         block/.style={rectangle, minimum size=1cm,draw,thick},
         operator/.style={circle,draw, minimum size=5mm, inner sep=0mm, thick},
         gain/.style={isosceles triangle,draw,isosceles triangle apex angle=60,thick}]
\node[name=input, coordinate] at (0,0) {};
\node[block, name=filtror, right of=input] {$\hat{H}(f)$};
\node[name=split,right of=filtror, xshift=-1cm, coordinate] {};

\node[block, name=delay, below right of=split] {Retardo $T_s$};
\node[gain, name=amp, right of=delay] {$2$};
\node[operator, name=mult, above right of=amp] {$\times$};
\node[block, name=int, right of=mult] {$\frac{1}{T_s}\int_{t-T_s}^t $};
\node[coordinate, name=sampler, right of=int] {};
\node[block, name=comp, right of=sampler] {Comp. $\pm$};

\draw[thick,->] (int.east)  -- ($(sampler) + (-.3,0)$) -- ($ (sampler) + (.3,.5) $) node[left, xshift=-5] {$t_k$};
\draw[thick,->] ($ (sampler) + (0.3,0) $) -- (comp.west);

\node[name=output, right of=comp] {$m_k$};

\draw[thick,->] (input)--(filtror);
\draw[thick,->] (filtror) -- (split) |- (delay.west);
\draw[thick,->] (delay.east) |- (amp.west);
\draw[thick,->] (amp.east) -| (mult.south);
\draw[thick,->] (split) -- (mult.west) node[midway] {$x_c(t)+n(t)$}; 
\draw[thick,->] (mult.east) -- (int.west);
\draw[thick,->] (comp.east) -- (output);

\end{tikzpicture}

\end{center}

\begin{enumerate}
\item Sea la onda BPSK $x_c(t) = \sum_k A p(t-kT_s)\cos(\omega_c t+a_k\pi),$
con $a_k\in\{0,1\}$, pulsos rectangulares NRZ,  y $\omega_c$
m\'ultiplo de $\omega_s$. Ignoramos por el momento el ruido. Sea $z_k$ el valor de la muestra en $t_k=(k+1)T_s$; el comparador da $m_k =1$ si $z_k>0$, y $m_k=0$ en otro caso. Expresar $m_k$ en función de $a_k$ y $a_{k-1}$.

\item Supongamos que se detecta la secuencia $m_1, m_2, \ldots, m_8 = 11000110$. Asumiendo que $a_0=1$, determinar cuál fue el mensaje $a_1, a_2, \ldots, a_8$ transmitido.

\item Consideramos ahora el ruido Gaussiano. Mediante un estudio que no se hará aquí, se puede obterner la fórmula
de probabilidad de error $ P_e = \frac{1}{2}\exp\left(-\frac{\bar{E}_p}{\eta}\right)$, donde $\bar{E}_p$ es la energía media del pulso transmitido, y
$\frac{\eta}{2}$ la densidad espectral de potencia del ruido blanco del canal. Tomando valores para $P_e = 10^{-i}, \ i=2,\ldots,10$, hacer una gráfica que compare los requerimientos de $\bar{E}_p/{\eta}$ con el caso del detector BPSK coherente. ¿Cuál es más eficiente?

\item Evaluar la potencia media requerida en el detector de arriba para una transmisión binaria de 100kbps y  un nivel de ruido $\frac{\eta}{2}=10^{-12}\text{V}^2/\text{Hz}$, si se desea $P_e=10^{-3}$.
\end{enumerate}

\end{ejercicio}

\begin{solucion}

        \begin{enumerate}
                \item La señal a la entrada del integrador (ignorando el ruido), denotando $\theta_k=a_k\pi$ es:
                \begin{align*}
                        y(t) &= 2x_c(t)x_c(t-T_s) \\
                        &= 2\left(\sum_k Ap(t-kT_s)\cos(\omega_c t + \theta_k)\right)\left(\sum_{k'} Ap(t-T_s-k'T_s)\cos(\omega_c t + \theta_k')\right)
                \end{align*}
                Notar que en el producto, los pulsos están desfasados y por lo tanto los productos cruzados dan $0$, excepto cuando $k'=k-1$, que da $p^2=1=p$, de donde:
                \begin{align*}
                        y(t) &= 2\sum_k Ap(t-kT_s)\cos(\omega_c t + \theta_k)\cos(\omega_c t + \theta_{k-1}) \\
                             &= \sum_k Ap(t-kT_s)\cos(\theta_k-\theta_{k-1}) + \ldots \\
                \end{align*}
                donde despreciamos los términos de alta frecuencia que se irán al filtrar (aquí usamos que en un tiempo de símbolo hay una cantidad entera de oscilaciones de la portadora, por hipótesis).

                Integrando en un período hallamos $z_k$:
                \begin{equation*}
                        z_k = \frac{1}{T_s}\int_{kT_s}^{(k+1)T_s} y(t)dt = \frac{1}{T_s}\int_{kT_s}^{(k+1)T_s} Ap(t-kT_s)\cos(\theta_k-\theta_{k-1}) = A\cos(\theta_k-\theta_{k-1})    
                \end{equation*}
                Notar que $z_k = 2A$ si $\theta_k \neq \theta_{k-1}$ y $0$ si $\theta_k=\theta_{k-1}$. Por eso se denomina demodulador diferencial, solo detecta cambios de símbolo o diferencias de fase.
                Resumiendo:
                \begin{equation*}
                        \begin{cases}
                                m_k = 1 & \text{si }a_k\neq a_{k-1}\\
                                m_k = 0 & \text{si }a_k = a_{k-1}
                        \end{cases}
                \end{equation*}
                
                \item Usando la idea anterior, $a_1\ldots a_8 = 01111011$.
                
                \item Tomando $\gamma = \bar{E}_p/\eta$, se tiene que:
                 \begin{align*}
                        P_e(\text{BPSK coh}) &= Q(\sqrt{2\gamma}), \\
                        P_e(\text{DBPSK} ) &= \frac{1}{2} e^{-\gamma}.
                 \end{align*}
                
                 Usando la aproximación $Q(x) \approx \frac{1}{\sqrt{2\pi} x}e^{-x^2/2}$, podemos escribir:
                 \begin{equation*}
                        \frac{P_e(\text{DBPSK})}{P_e(\text{BPSK coh})} = \frac{\frac{1}{2}e^{-\gamma}}{\frac{1}{\sqrt{4\pi\gamma}}e^{-\gamma}} = \sqrt{\pi\gamma}.
                 \end{equation*}
                 
                 La performance de DBPSK es levemente peor (en escala logarítmica). Graficamos a continuación ambas en función de $\gamma$:
                 \begin{center}
                        \begin{tikzpicture}

\begin{axis}[   width = 0.8\textwidth, height = 0.5\textwidth,
                xlabel=$\gamma$ (dB),
                ylabel=$P_e$,
                x label style={anchor=north west},
                y label style={anchor=south east},
                xmin=3,xmax=15,ymin=-18,ymax=0, grid,
                xtick = {3,6,9,12,15},
                ytick = {-15,-10,-5,0},
                yticklabels = {$10^{-15}$,$10^{-10}$,$10^{-5}$,$10^0$},
            ]
    \addplot[color=teal, very thick, mark=none, solid]
        table[row sep={\\}]
        {
            \\
            3.0  -1.1675614090044737  \\
            3.121212121212122  -1.1920870718343202  \\
            3.2424242424242427  -1.2173068908459077  \\
            3.3636363636363638  -1.24324051292128  \\
            3.4848484848484853  -1.2699081410132667  \\
            3.606060606060606  -1.2973305498840992  \\
            3.7272727272727275  -1.3255291022894815  \\
            3.848484848484848  -1.354525765620719  \\
            3.9696969696969697  -1.3843431290178796  \\
            4.090909090909091  -1.4150044209673063  \\
            4.212121212121211  -1.446533527397203  \\
            4.333333333333334  -1.4789550102853808  \\
            4.454545454545454  -1.5122941267936645  \\
            4.575757575757576  -1.546576848943866  \\
            4.696969696969697  -1.5818298838506517  \\
            4.818181818181818  -1.6180806945270654  \\
            4.9393939393939394  -1.6553575212789158  \\
            5.0606060606060606  -1.693689403704695  \\
            5.181818181818182  -1.7331062033181661  \\
            5.303030303030303  -1.7736386268112432  \\
            5.424242424242425  -1.815318249975289  \\
            5.545454545454546  -1.8581775422994584  \\
            5.666666666666666  -1.9022498922652615  \\
            5.787878787878787  -1.9475696333570405  \\
            5.909090909090909  -1.9941720708086301  \\
            6.03030303030303  -2.042093509107033  \\
            6.151515151515151  -2.0913712802745428  \\
            6.272727272727273  -2.142043772951343  \\
            6.3939393939393945  -2.19415046230123  \\
            6.515151515151515  -2.2477319407637797  \\
            6.636363636363637  -2.30282994967689  \\
            6.757575757575758  -2.3594874117943436  \\
            6.878787878787879  -2.4177484647237337  \\
            7.0  -2.4776584953107834  \\
            7.121212121212121  -2.539264174996854  \\
            7.242424242424242  -2.6026134961771885  \\
            7.363636363636363  -2.667755809588213  \\
            7.484848484848485  -2.7347418627530184  \\
            7.606060606060606  -2.8036238395149704  \\
            7.727272727272727  -2.8744554006902616  \\
            7.848484848484848  -2.947291725871049  \\
            7.96969696969697  -3.022189556411767  \\
            8.09090909090909  -3.0992072396320856  \\
            8.212121212121211  -3.1784047742709647  \\
            8.333333333333334  -3.259843857227199  \\
            8.454545454545455  -3.343587931622874  \\
            8.575757575757576  -3.429702236227177  \\
            8.696969696969697  -3.5182538562790664  \\
            8.818181818181818  -3.609311775748372  \\
            8.93939393939394  -3.7029469310760783  \\
            9.06060606060606  -3.7992322664356135  \\
            9.181818181818182  -3.898242790558224  \\
            9.303030303030303  -4.000055635166685  \\
            9.424242424242424  -4.104750115062883  \\
            9.545454545454547  -4.2124077899160675  \\
            9.666666666666666  -4.3231125277999105  \\
            9.787878787878787  -4.436950570527878  \\
            9.90909090909091  -4.554010600837798  \\
            10.03030303030303  -4.674383811477979  \\
            10.151515151515152  -4.798163976248683  \\
            10.272727272727273  -4.925447523054305  \\
            10.393939393939393  -5.05633360902317  \\
            10.515151515151516  -5.190924197753473  \\
            10.636363636363637  -5.329324138745508  \\
            10.757575757575758  -5.471641249082107  \\
            10.878787878787879  -5.617986397420906  \\
            11.0  -5.768473590363841  \\
            11.121212121212121  -5.923220061271204  \\
            11.242424242424242  -6.082346361589429  \\
            11.363636363636365  -6.245976454763728  \\
            11.484848484848484  -6.4142378128087785  \\
            11.606060606060606  -6.587261515612687  \\
            11.727272727272727  -6.765182353051539  \\
            11.848484848484848  -6.9481389299941805  \\
            11.96969696969697  -7.136273774278913  \\
            12.09090909090909  -7.329733447746334  \\
            12.212121212121211  -7.528668660414745  \\
            12.333333333333334  -7.733234387887093  \\
            12.454545454545453  -7.943589992080915  \\
            12.575757575757576  -8.159899345375353  \\
            12.696969696969697  -8.382330958271883  \\
            12.818181818181817  -8.61105811066832  \\
            12.93939393939394  -8.846258986848282  \\
            13.060606060606062  -9.088116814291249  \\
            13.181818181818182  -9.336820006411482  \\
            13.303030303030303  -9.592562309336872  \\
            13.424242424242426  -9.855542952842086  \\
            13.545454545454545  -10.125966805553666  \\
            13.666666666666668  -10.404044534547918  \\
            13.787878787878789  -10.689992769465906  \\
            13.909090909090908  -10.984034271273494  \\
            14.030303030303031  -11.286398105797813  \\
            14.15151515151515  -11.597319822175358  \\
            14.272727272727273  -11.917041636350815  \\
            14.393939393939394  -12.24581261976941  \\
            14.515151515151514  -12.583888893409938  \\
            14.636363636363637  -12.931533827309556  \\
            14.757575757575758  -13.289018245735717  \\
            14.878787878787879  -13.656620638165196  \\
            15.0  -14.034627376234518  \\
        }
        ;
    \addlegendentry {DBPSK}
    \addplot[color=red, very thick, mark=none, solid]
        table[row sep={\\}]
        {
            \\
            3.0  -1.6405742076544345  \\
            3.121212121212122  -1.669633827680205  \\
            3.2424242424242427  -1.6994119090038986  \\
            3.3636363636363638  -1.7299279160069392  \\
            3.4848484848484853  -1.761201865445318  \\
            3.606060606060606  -1.7932543422837437  \\
            3.7272727272727275  -1.826106515978037  \\
            3.848484848484848  -1.8597801572182753  \\
            3.9696969696969697  -1.894297655145531  \\
            4.090909090909091  -1.9296820350554105  \\
            4.212121212121211  -1.9659569766019847  \\
            4.333333333333334  -2.003146832516071  \\
            4.454545454545454  -2.0412766478522313  \\
            4.575757575757576  -2.0803721797792574  \\
            4.696969696969697  -2.120459917929319  \\
            4.818181818181818  -2.161567105321405  \\
            4.9393939393939394  -2.2037217598751084  \\
            5.0606060606060606  -2.2469526965312725  \\
            5.181818181818182  -2.2912895499964896  \\
            5.303030303030303  -2.3367627981289116  \\
            5.424242424242425  -2.3834037859833486  \\
            5.545454545454546  -2.431244750534125  \\
            5.666666666666666  -2.4803188460947014  \\
            5.787878787878787  -2.5306601704536096  \\
            5.909090909090909  -2.582303791746801  \\
            6.03030303030303  -2.6352857760870916  \\
            6.151515151515151  -2.689643215971962  \\
            6.272727272727273  -2.7454142594916084  \\
            6.3939393939393945  -2.8026381403597127  \\
            6.515151515151515  -2.8613552087901146  \\
            6.636363636363637  -2.92160696324316  \\
            6.757575757575758  -2.983436083066217  \\
            6.878787878787879  -3.0468864620535743  \\
            7.0  -3.112003242951578  \\
            7.121212121212121  -3.1788328529357055  \\
            7.242424242424242  -3.2474230400869377  \\
            7.363636363636363  -3.317822910895672  \\
            7.484848484848485  -3.3900829688221266  \\
            7.606060606060606  -3.4642551539430917  \\
            7.727272727272727  -3.540392883715711  \\
            7.848484848484848  -3.618551094889824  \\
            7.96969696969697  -3.6987862866013623  \\
            8.09090909090909  -3.781156564680149  \\
            8.212121212121211  -3.8657216872064755  \\
            8.333333333333334  -3.952543111351749  \\
            8.454545454545455  -4.041684041539546  \\
            8.575757575757576  -4.133209478964443  \\
            8.696969696969697  -4.227186272507052  \\
            8.818181818181818  -4.323683171084765  \\
            8.93939393939394  -4.422770877478903  \\
            9.06060606060606  -4.524522103680035  \\
            9.181818181818182  -4.629011627794504  \\
            9.303030303030303  -4.736316352556355  \\
            9.424242424242424  -4.846515365490199  \\
            9.545454545454547  -4.959690000771722  \\
            9.666666666666666  -5.075923902834024  \\
            9.787878787878787  -5.195303091769213  \\
            9.90909090909091  -5.31791603057617  \\
            10.03030303030303  -5.443853694306821  \\
            10.151515151515152  -5.573209641164683  \\
            10.272727272727273  -5.706080085611124  \\
            10.393939393939393  -5.842563973536148  \\
            10.515151515151516  -5.982763059552364  \\
            10.636363636363637  -6.126781986472189  \\
            10.757575757575758  -6.274728367030339  \\
            10.878787878787879  -6.426712867915166  \\
            11.0  -6.582849296174322  \\
            11.121212121212121  -6.743254688062123  \\
            11.242424242424242  -6.908049400397799  \\
            11.363636363636365  -7.07735720450581  \\
            11.484848484848484  -7.251305382811475  \\
            11.606060606060606  -7.430024828167186  \\
            11.727272727272727  -7.613650145986593  \\
            11.848484848484848  -7.802319759266413  \\
            11.96969696969697  -7.996176016577684  \\
            12.09090909090909  -8.195365303110657  \\
            12.212121212121211  -8.400038154859919  \\
            12.333333333333334  -8.61034937603868  \\
            12.454545454545453  -8.826458159813845  \\
            12.575757575757576  -9.048528212455938  \\
            12.696969696969697  -9.276727881000662  \\
            12.818181818181817  -9.511230284521636  \\
            12.93939393939394  -9.75221344911665  \\
            13.060606060606062  -9.999860446712638  \\
            13.181818181818182  -10.254359537797605  \\
            13.303030303030303  -10.515904318190808  \\
            13.424242424242426  -10.784693869965505  \\
            13.545454545454545  -11.060932916642045  \\
            13.666666666666668  -11.344831982772224  \\
            13.787878787878789  -11.636607558039294  \\
            13.909090909090908  -11.93648226600157  \\
            14.030303030303031  -12.244685037611198  \\
            14.15151515151515  -12.56145128964326  \\
            14.272727272727273  -12.887023108174441  \\
            14.393939393939394  -13.221649437254078  \\
            14.515151515151514  -13.565586272914876  \\
            14.636363636363637  -13.919096862674358  \\
            14.757575757575758  -14.282451910682568  \\
            14.878787878787879  -14.655929788675973  \\
            15.0  -15.039816752902032  \\
        }
        ;
    \addlegendentry {Coherente}
\end{axis}
\end{tikzpicture}

                 \end{center}

        \end{enumerate}
        
\end{solucion}


\begin{ejercicio}
Un broadcaster tiene un flujo de video de 30 Mbps que se desea transportar via satélite. Se decide utilizar modulación 4-QAM con pulsos de Nyquist de coseno elevado con exceso de ancho de banda $\alpha=0.2$.
\begin{enumerate}
\item Indicar el ancho de banda $B$ necesario en el canal pasabanda para este transporte.
\item Bosquejar la densidad espectral de potencia de la señal modulada.
\item Proponer un detector coherente para estos pulsos que no introduzca ISI, y deducir la fórmula para la probabilidad de error de detección en base a los parámetros de modulación, en particular la amplitud $A$ de los pulsos recibidos y la densidad espectral de potencia del ruido blanco en el canal. Evaluar para el caso $A=3V$ y $\frac{\eta}{2} = 10^{-8}\unit{\square\volt/\hertz}$.
\item Luego de realizados los cálculos, el proveedor de espacio satelital avisa que no podrá ofrecer el ancho de banda $B$, pero si se podrá ofrecer dos espacios de ancho de banda $\frac{B}{2}$. ¿Es posible transmitir los 30 Mbps entre esos dos canales mediante modulación 4QAM en c/u, con el mismo rolloff en
los pulsos de Nyquist? Si es así indicar cómo hacerlo.
\end{enumerate}
\end{ejercicio}

\begin{solucion}

        \begin{enumerate}
                \item Necesita enviar $30$ Mbps que corresponden a $15$ MBaudios a $2$ bits por símbolo. La frecuencia de símbolo es entonces $15\unit{\mega\hertz}$.
                
                El espectro equivalente en banda base es $S_{x_{ce}} = E[c_k^2]f_s |\hat{P}(f)|^2$, donde en este caso $\hat{P}$ es un pulso de Nyquist con exceso $\alpha=0.2$. Éste tiene un ancho de banda $(1+\alpha)f_s/2$, y al modularlo requiere $(1+\alpha)f_s$ de ancho de banda.

                Por lo tanto, $B=(1.2)f_s = 18 \unit{\mega\hertz}$.

                \item Observando que $E[c_k^2] = 2 A^2$ en $4-$QAM, el espectro equivalente pasabajos tiene ancho de banda $\pm (1+\alpha)f_s/2$ y alto $2A^2 T_s$. Al modularlo la altura queda $A^2 T_s/2$ en cada banda.
                
                Bosquejo (para $f>0$):
                \begin{center}
                        \begin{tikzpicture}
                                \begin{axis}[x=1cm, y=1cm, axis lines=middle,  
                                        xlabel=$f$,
                                        ylabel=$S_{x_c}(f)$,,
                                        x label style={anchor=north west},
                                        y label style={anchor=south east},
                                        xmin=-.1,xmax=8,ymin=-.1,ymax=1.5,
                                        xtick={3.6,5,6.4}, ytick=\empty,
                                        xticklabels={\footnotesize $f_c-\frac{(1+\alpha)f_s}{2}$,\footnotesize $f_c$,\footnotesize  $f_c+\frac{(1+\alpha)f_s}{2}$},
                                        clip=false,
                                    ]
                                \addplot+[solid, mark=none, very thick, teal, domain = (4:6)]  {1} node[midway, above right] {$\frac{1}{2}A^2 T_s$};
                                \addplot+[solid, mark=none, very thick, teal, domain = (3.6:4), samples=100]  {.5+.5*cos(2*pi*(x-4)/0.8 r)};
                                \addplot+[solid, mark=none, very thick, teal, domain = (6:6.4), samples=100]  {.5+.5*cos(2*pi*(x-6)/0.8 r)};
                                \draw[dashed] (axis cs:5,0) -- (axis cs:5,1.5);
                                \end{axis}
                        \end{tikzpicture}
                \end{center}

                \item Se desea un detector que no introduzca ISI, por lo que debe dejarse pasar el pulso completo. El detector es:
                
                \begin{center}
                \begin{tikzpicture}
                        [node distance=1.5cm,>=stealth, auto,
                        block/.style={rectangle, minimum size=1cm,draw,thick},
                        operator/.style={circle,draw, minimum size=5mm, inner sep=0mm, thick},
                        gain/.style={isosceles triangle,draw,isosceles triangle apex angle=60,thick}]

                \node[anchor=west] (input) at (0,0) {$x_c(t)+n(t)$};
                \node[coordinate, right of=input] (split) {};
                \node[operator,  right of=split] (mult1) {$\times$};
                \node[below of=mult1] (cos) {$2\cos(\omega_ct)$};
                \node[operator, below of=cos] (mult2) {$\times$};
                \node[below of=mult2] (sin) {$-2\sin(\omega_ct)$};

                \draw[thick,->] (input) -- (mult1);
                \draw[thick,->] (split) |- (mult2);
                \draw[thick,->] (cos) -- (mult1);
                \draw[thick,->] (sin) -- (mult2);

                \node[block, name=filtroi, right of=mult1] {$\hat{H}(f)$};
                \node[block, name=filtroq, right of=mult2] {$\hat{H}(f)$};
                \draw[thick,->] (mult1) -- (filtroi);
                \draw[thick,->] (mult2) -- (filtroq);
                
                \node[coordinate, name=sampleri, right of=filtroi] {};
                \node[block, name=compi, right of=sampleri, xshift=.5cm] {Comp. $\pm$};
                \draw[thick,->] (filtroi.east)  -- ($(sampleri) + (-.3,0)$) -- ($ (sampleri) + (.3,.5) $) node[left, xshift=-5] {$t_k$};
                \draw[thick,->] ($ (sampleri) + (0.3,0) $) -- (compi.west);
                \node[name=ak, right of = compi, xshift=1cm] (ak) {$a_k$};
                \draw[thick,->] (compi.east) -- (ak);

                \node[coordinate, name=samplerq, right of=filtroq] {};
                \node[block, name=compq, right of=samplerq, xshift=.5cm] {Comp. $\pm$};
                \draw[thick,->] (filtroq.east)  -- ($(samplerq) + (-.3,0)$) -- ($ (samplerq) + (.3,.5) $) node[left, xshift=-5] {$t_k$};
                \draw[thick,->] ($ (samplerq) + (0.3,0) $) -- (compq.west);
                \node[name=ak, right of = compq, xshift=1cm] (bk) {$b_k$};
                \draw[thick,->] (compq.east) -- (bk);

                \end{tikzpicture}
                \end{center}

                El filtro $\hat{H}(f)$ en recepción es un pasabajos (ideal) de ancho $(1+\alpha)f_s/2$.

                La señal demodulada a la entrada del comparador está intacta, por lo que tendrá amplitud $\pm A$. Queda evaluar la potencia del ruido. La densidad espectral del ruido blanco en banda base es $\eta$ y al pasar por el pasabajos ideal se tiene $\sigma^2 = \eta (1+\alpha)f_s$. La probabilidad de error queda:
                \begin{equation*}
                        P_e(a_k) = P_e(b_k) = Q\left(\sqrt{\frac{A^2}{\eta (1+\alpha)f_s}}\right) \Rightarrow P_e \approx 2Q\left(\sqrt{\frac{A^2}{\eta (1+\alpha)f_s}}\right)
                \end{equation*}

                Para los valores anteriores, $P_e = 2Q(5) = 5.7 \times 10^{-7}$.

                \item Es posible, dividiendo el flujo en dos canales de $15$ Mbps, usando $f_s=7.5 \unit{\mega\hertz}$, usando el exceso $\alpha=0.2$, el canal queda de $B/2=9\unit{\mega\hertz}$ que es lo ofrecido por el operador. El problema es que al dividir, la probabilidad de error aumenta si se mantiene la potencia constante, ya que $f_s$ es la mitad. Dicho de otro modo, se debe aumentar la potencia en $3$ dB.
        \end{enumerate}
\end{solucion}


\begin{ejercicio}
Sea un sistema de transmisión digital pasabanda, en donde envía la señal con envolvente compleja
\begin{equation*}
x_{ce}(t)=\sum z_k p_{[0,T_s]}(t),
\end{equation*}
en donde los símbolos complejos $z_k$ son independientes y equiprobables. Se utiliza la siguiente constelación de símbolos

\begin{figure}[h!]
        \centering
\begin{tikzpicture}[scale=1]
                \draw [->] (-3,0) -- (3,0) node [anchor=north] {$I$};
                \draw [->] (0,-3) -- (0,3) node [anchor=west] {$Q$};
                \draw [teal,fill] (2,2) circle (0.14);
                \draw [teal,fill] (-2,2) circle (0.14);
                \draw [teal,fill] (2,-2) circle (0.14);
                \draw [teal,fill] (-2,-2) circle (0.14);
                \draw [teal,fill] (-1,2) circle (0.14);
                \draw [teal,fill] (1,-2) circle (0.14);
                \draw [teal,fill] (1,2) circle (0.14);
                \draw [teal,fill] (-1,-2) circle (0.14);
                \draw [teal,fill] (-2,1) circle (0.14);
                \draw [teal,fill] (2,-1) circle (0.14);
                \draw [teal,fill] (-2,-1) circle (0.14);
                \draw [teal,fill] (2,1) circle (0.14);
                \draw [teal,fill] (-1,1) circle (0.14);
                \draw [teal,fill] (1,-1) circle (0.14);
                \draw [teal,fill] (-1,-1) circle (0.14);
                \draw [teal,fill] (1,1) circle (0.14);
%      \node at(3,3) {16-QAM};
                \foreach \x/\y in {-2,-1/-,1/{},2}
                        {\draw [fill] (\x,0)  node [anchor=north] {\y$A$};
                \draw [fill] (0,\x)  node [anchor=east] {\y$A$};}
%
\end{tikzpicture}
\end{figure}

\begin{enumerate}
\item Buscar una codificación de Gray adecuada para el diagrama, hallar la densidad espectral de potencia en banda base equivalente $S_{ce}(f)$ y la energía por bit $\bar{E_b}$.
\item Diseñar un diagrama de bloques de un demodulador especificando los umbrales y los filtros utilizados.
\item Expresar la probabilidad de error en función de A y $\sigma$ para $a_k$ y $b_k $ y dar una aproximación para el error de símbolo. Suponer que se tiene ruido blanco Gaussiano aditivo con varianza $\sigma^2$.
\item Para $A=6$ y $\sigma=1$, evaluar la probabilidad de error de símbolo y compararla con la probabilidad que se obtiene si se desprecian los errores que se escapan del cuadrante del símbolo.
\item Expresar la probabilidad anterior en función de $\bar{E_p}$ y $\eta$, donde $\frac{\eta}{2}$ es la densidad espectral del ruido.
\item Suponiendo que $a_k$ y $b_k$ tienen sus valores en $(2A,\alpha A, -\alpha A, 2A)$, hallar el $\alpha$ que hace óptimo la probabilidad de error de símbolo, especificando cual es esta probabilidad.
\end{enumerate}
\end{ejercicio}

\begin{solucion}
        \begin{enumerate}
                \item Codificación de Gray:
                \begin{center}
                        \begin{tikzpicture}
                                \begin{axis}[x=1.5cm, y=1.5cm, axis lines=middle,  
                                        xlabel=$a_k$,
                                        ylabel=$b_k$,
                                        x label style={anchor=north west},
                                        y label style={anchor=south east},
                                        xmin=-2.5,xmax=2.5,ymin=-2.5,ymax=2.5,
                                        xtick={-2,-1,0,1,2}, ytick={-2,-1,0,1,2},
                                        xticklabels={$-2A$, $-A$, $0$, $A$, $2A$},
                                        yticklabels={$-2A$,$-A$, $0$, $A$, $2A$},
                                        clip=false,
                                    ]
                                \addplot+[only marks, teal, mark options={fill=teal}] coordinates {
                                        (-2,-1) (-2,1) (-2,-2) (-2,2)
                                        (-1,-1) (-1,1) (-1,-2) (-1,2)
                                        (1,-1) (1,1) (1,-2) (1,2)
                                        (2,-1) (2,1) (2,-2) (2,2)
                                };
                                \node[below] at (axis cs:-2,-2) { (0000)};
                                \node[below] at (axis cs:-2,-1) { (0100)};
                                \node[below] at (axis cs:-1,-2) { (0001)};
                                \node[below] at (axis cs:-1,-1) { (0101)};
                                \node[below] at (axis cs:1,-2) { (0011)};
                                \node[below] at (axis cs:1,-1) { (0111)};
                                \node[below] at (axis cs:2,-2) { (0010)};
                                \node[below] at (axis cs:2,-1) { (0110)};
                                \node[below] at (axis cs:-2,2) { (1000)};
                                \node[below] at (axis cs:-2,1) { (1100)};
                                \node[below] at (axis cs:-1,2) { (1001)};
                                \node[below] at (axis cs:-1,1) { (1101)};
                                \node[below] at (axis cs:1,2) { (1011)};
                                \node[below] at (axis cs:1,1) { (1111)};
                                \node[below] at (axis cs:2,2) { (1010)};
                                \node[below] at (axis cs:2,1) { (1110)};

                                \end{axis}
                        \end{tikzpicture}
                \end{center}

                Para la densidad espectral, observemos que la constelación tiene $\mu=0$ y necesitamos $E[c_k^2]$. Por simetría, lo calculamos en un solo cuadrante, obteniendo $E[c_k^2] = 2A^2(1/4) + 5A^2(1/2) + 8A^2(1/4) = 5A^2$. Por lo tanto:
                \begin{equation*}
                        S_{x_{ce}}(f) = 5A^2 T_s \sinc^2(\pi T_s f),
                \end{equation*}
                donde se usó la transformada del pulso NRZ.

                La energía media del pulso es $\bar{E}_p = E[c_k^2]||p||^2 = 5A^2T_s$ y la energía por bit es $\bar{E}_b = \bar{E}_p/4 = (5/4)A^2 T_s$.

                \item Se hace detección coherente como en $16$ QAM con filtro acoplado. El único detalle es que los umbrales son $-(3/2)A$, $0$ y $(3/2)A$ porque los símbolos no están equidistantes.
                
                \item Calculamos $P_e(a_k)=P_e(b_k)$ en función de $A/\sigma$ observando las distancias. Queda:
                 \begin{align*}
                        P_e(a_k) &= \frac{1}{4}\left[Q(A/(2\sigma)) + Q(A/(2\sigma)) + Q(A/\sigma) + Q(A/\sigma) + Q(A/(2\sigma))+ Q(A/(2\sigma))\right] \\
                        &=Q(A/(2\sigma))+ \frac{1}{2}Q(A/\sigma)
                 \end{align*}
                 Por lo tanto, la probabilidad de error total es:
                 \begin{equation*}
                        P_e \approx P_e(a_k)+P_e(b_k) = 2Q(A/(2\sigma))+ Q(A/\sigma)
                 \end{equation*}
                 
                 \item Para $(A/\sigma) = 6$ queda $P_e = 0.0026$. El término $Q(A/\sigma)$ que corresponde a cambiarse de cuadrante es despreciable frente al otro. Los errores se dan solo dentro del cuadrante. $P_e \approx 2Q(A/(2\sigma))$.

                 \item Recordando que $\bar{E}_p = 5A^2T_s$ y $\sigma^2=\eta f_s$, se tiene que:
                  \begin{equation*}
                        P_e = 2Q\left(\sqrt{\frac{\bar{E}_p}{5\eta}}\right) = 2Q\left(\sqrt{\frac{4\bar{E}_b}{5\eta}}\right)
                  \end{equation*}

                  \item Lo mejor es poner los símbolos equidistantes, es decir:
                   \begin{equation*}
                        (2A-\alpha A) = \alpha A - (-\alpha)A,
                   \end{equation*}
                   de donde $\alpha=2/3$. En ese caso, la constelación queda la de QAM-16, pero cambiando $A$ por $2/3 A$, así que $P_e \approx 3Q\left(2A/(3\sigma)\right)$.
        \end{enumerate}
\end{solucion}
\end{document}